% Luc Destombe --- licence CC BY-NC-SA 4.0
% https://creativecommons.org/licenses/by-nc-sa/4.0/deed.fr
% Fichier autonome : compiler avec pdflatex, deux fois (signets, nombre de pages).
\documentclass[11pt,a4paper]{article}
\makeatletter
\newif\iflycee@unecolonne
\lycee@unecolonnetrue
\newif\iflycee@palatino
\lycee@palatinotrue

\RequirePackage[utf8]{inputenc}
\RequirePackage[T1]{fontenc}
\RequirePackage[french]{babel}
\RequirePackage{etoolbox}

\newcommand{\ShorthandsOff}{\@ifundefined{shorthandoff}{}{\shorthandoff{;:!?}}}
\newcommand{\ShorthandsOn}{\@ifundefined{shorthandon}{}{\shorthandon{;:!?}}}
\ShorthandsOff
\AtBeginDocument{\ShorthandsOn}
\AtBeginEnvironment{tikzpicture}{\ShorthandsOff}

\RequirePackage{geometry}
\RequirePackage{fancyhdr}
\RequirePackage{lastpage}
\RequirePackage{multicol}
\RequirePackage{array}
\RequirePackage{enumitem}
\RequirePackage{xcolor}
\RequirePackage{needspace}

\RequirePackage{amsmath,amssymb}
\RequirePackage{mathpazo}
\DeclareMathSizes{10.95}{10.95}{8}{5.5}
\begingroup\catcode`\_=\active \catcode`\^=\active
\gdef_#1{\sb{\scriptscriptstyle #1}}
\gdef^#1{\sp{\scriptscriptstyle\lycee@moinsepais #1}}
\endgroup
\newcommand\lycee@moins{\mathbin{%
  \setbox\z@\hbox{$\m@th\scriptscriptstyle\lycee@moinsorig$}%
  \dimen@=.55\wd\z@ \advance\dimen@-.5pt
  \hbox to.75\wd\z@{\hss$\m@th\scriptscriptstyle\vcenter{\hbox{%
    \tikz\draw[line width=.5pt,line cap=round](0,0)--(\the\dimen@,0);}}$\hss}}}
\begingroup\lccode`\~=`\- \lowercase{\endgroup\let~}\lycee@moins
\newcommand\lycee@moinsepais{\mathcode`\-="8000 }
\AtBeginDocument{\mathchardef\lycee@moinsorig=\mathcode`\-\relax}
\AtBeginDocument{\catcode`\_=12 \mathcode`\_="8000
  \catcode`\^=12 \mathcode`\^="8000 }
\newcommand\lycee@exposants{%
  \fontdimen13\textfont2=.48\fontdimen6\textfont2
  \fontdimen14\textfont2=.48\fontdimen6\textfont2
  \fontdimen15\textfont2=.48\fontdimen6\textfont2
  \fontdimen13\scriptfont2=.48\fontdimen6\scriptfont2
  \fontdimen14\scriptfont2=.48\fontdimen6\scriptfont2
  \fontdimen15\scriptfont2=.48\fontdimen6\scriptfont2}
\every@math@size\expandafter{\the\every@math@size\lycee@exposants}
\newlength{\retraitcalculs}
\setlength{\retraitcalculs}{2em}

\RequirePackage{tikz}
\usetikzlibrary{arrows.meta,calc,babel,shapes.misc}
\RequirePackage[most]{tcolorbox}
\tcbuselibrary{skins,breakable}

\definecolor{coulDef}{HTML}{1F4E79}
\definecolor{coulProp}{HTML}{7030A0}
\definecolor{coulRem}{HTML}{C55A11}
\definecolor{coulEx}{HTML}{2E7D32}
\definecolor{coulExo}{HTML}{B03A2E}
\definecolor{coulPart}{HTML}{1F4E79}
\definecolor{coulMeth}{HTML}{1B6E4F}
\definecolor{coulCourbe}{HTML}{0B5ED7}
\definecolor{coulCourbeB}{HTML}{C00000}
\definecolor{coulCourbeC}{HTML}{2E7D32}
\definecolor{coulGrille}{HTML}{8C8C8C}
\definecolor{coulRep}{HTML}{1B6E4F}
\definecolor{coulBar}{HTML}{2E7D32}

\newcommand{\R}{\mathbb{R}}
\newcommand{\Rs}{\mathbb{R}^{*}}
\renewcommand{\le}{\leqslant}
\renewcommand{\ge}{\geqslant}
\renewcommand{\approx}{\simeq}
\newcommand{\vect}[1]{\overrightarrow{#1}}
\newcommand{\norme}[1]{\left\lVert #1 \right\rVert}
\newcommand{\intervalle}[2]{\left[#1\,;#2\right]}
\RequirePackage{cancel}
\renewcommand{\CancelColor}{\color{coulExo}}
\newcommand{\fois}[1]{\times{\color{coulExo}#1}}
\newcommand{\barre}[1]{\cancel{#1}}

\tikzset{grille/.style={coulGrille,line width=0.4pt}}
\newcommand{\quadrillage}[4]{\draw[grille,step=1] (#1,#2) grid (#3,#4);}
\tikzset{
  croix/.style={cross out,draw,line width=0.6pt,inner sep=0pt,minimum size=4pt},
  croixdroite/.style={croix,rotate=45,line width=0.8pt},
}
\newcommand{\pt}[3]{\path (#1) node[croix]{} node[#2,font=\small,inner sep=2.5pt]{$#3$};}

\newcommand{\lycee@repere}[4]{%
  \draw[grille,step=1] (#1,#2) grid (#3,#4);
  \draw[-{Stealth[length=2mm]},thick] (#1,0)--({#3+0.4},0) node[below left=-1pt]{$x$};
  \draw[-{Stealth[length=2mm]},thick] (0,#2)--(0,{#4+0.4}) node[below left=-1pt]{$y$};
  \node[below left,font=\scriptsize,inner sep=1.5pt] at (0,0) {$0$};}
\newcommand{\lycee@gradx}[1]{\foreach \k/\l in {#1}{%
  \draw (\k,0.07)--(\k,-0.07) node[below,font=\scriptsize,inner sep=1.5pt]{$\l$};}}
\newcommand{\lycee@grady}[1]{\foreach \k/\l in {#1}{%
  \draw (0.07,\k)--(-0.07,\k) node[left,font=\scriptsize,inner sep=1.5pt]{$\l$};}}
\newcommand{\lycee@intff}[2]{\left[#1\,;#2\right]}
\newcommand{\lycee@intfo}[2]{\left[#1\,;#2\right[}
\newcommand{\lycee@intof}[2]{\left]#1\,;#2\right]}
\newcommand{\lycee@intoo}[2]{\left]#1\,;#2\right[}
\newcommand{\lycee@ens}[1]{\left\{#1\right\}}
\AtBeginDocument{%
  \providecommand{\repere}{\lycee@repere}%
  \providecommand{\gradx}{\lycee@gradx}%
  \providecommand{\grady}{\lycee@grady}%
  \providecommand{\Cf}{\mathcal{C}\sb{\scriptscriptstyle f}}%
  \providecommand{\Cg}{\mathcal{C}\sb{\scriptscriptstyle g}}%
  \providecommand{\intff}{\lycee@intff}%
  \providecommand{\intfo}{\lycee@intfo}%
  \providecommand{\intof}{\lycee@intof}%
  \providecommand{\intoo}{\lycee@intoo}%
  \providecommand{\ens}{\lycee@ens}}

\newlist{questions}{enumerate}{3}
\setlist[questions,1]{label=\arabic*\textdegree),leftmargin=*,itemsep=2pt,topsep=3pt}
\setlist[questions,2]{label=\alph*),leftmargin=*,itemsep=1pt,topsep=2pt}
\setlist[questions,3]{label=\roman*),leftmargin=*,itemsep=1pt,topsep=1pt}
\setlist[itemize]{label=\textbullet,leftmargin=*,itemsep=2pt,topsep=3pt}

\newcommand{\besoinplace}[1]{\par\penalty\@M\begingroup
  \dimen@=#1\relax\dimen@ii=\pagegoal\advance\dimen@ii-\pagetotal
  \ifdim\dimen@>\dimen@ii\ifdim\dimen@ii>\z@\vfil\fi\break\fi\endgroup}

\newcommand{\lycee@pied}{\thepage/\pageref{LastPage}}
\newcommand{\entete}[2]{%
  \pagestyle{fancy}\fancyhf{}%
  \fancyhead[L]{\footnotesize\color{coulPart}\textbf{#1}}%
  \fancyhead[R]{\footnotesize\color{coulPart}\textbf{#2}}%
  \fancyfoot[C]{\footnotesize\lycee@pied}%
  \renewcommand{\headrulewidth}{0.4pt}%
  \renewcommand{\headrule}{\hbox to\headwidth{%
    \color{coulPart!40}\leaders\hrule height \headrulewidth\hfill}}}

\RequirePackage{ccicons}
\newcommand{\lycee@licence}{{\scriptsize\color{gray}%
  \copyright\ Luc Destombe --- \ccbyncsa\ CC BY-NC-SA 4.0}}
\newif\iflycee@licence \lycee@licencetrue
\newcommand{\sanslicence}{\lycee@licencefalse}
\AtBeginDocument{\iflycee@licence\AddToHookNext{shipout/foreground}{%
  \put(0,0){\rlap{\kern\dimexpr1in+\hoffset+\oddsidemargin+\textwidth\relax
    \llap{\raisebox{-\dimexpr1in+\voffset+\topmargin+\headheight+\headsep
      +\textheight+\footskip\relax}{\lycee@licence}}}}}\fi}

\newcommand{\formulesserrees}{%
  \setlength{\abovedisplayskip}{3pt}\setlength{\belowdisplayskip}{3pt}%
  \setlength{\abovedisplayshortskip}{2pt}\setlength{\belowdisplayshortskip}{3pt}}

\setlength{\parindent}{0pt}

\RequirePackage[implicit=false,hidelinks,pdfencoding=auto,psdextra,bookmarksopen,
  bookmarksopenlevel=1,pdfcreator={},pdfproducer={}]{hyperref}
\RequirePackage{bookmark}
\hypersetup{pdfauthor={Luc Destombe},pdfkeywords={CC BY-NC-SA 4.0}}
\newcounter{lycee@signet}
\newcommand{\signet}[3][]{\stepcounter{lycee@signet}%
  \edef\lycee@dest{\if\relax\detokenize{#1}\relax
    signet.\arabic{lycee@signet}\else#1\fi}%
  \hypertarget{\lycee@dest}{}\bookmark[level=#2,dest=\lycee@dest]{#3}}
\pdfstringdefDisableCommands{%
  \def\R{R}\def\vect#1{#1}\def\Cf{Cf}\def\Cg{Cg}\def\fois#1{×#1}%
  \def\barre#1{#1}\def\intervalle#1#2{[#1;#2]}\def\intff#1#2{[#1;#2]}%
  \def\textsuperscript#1{#1}\def\dfrac#1#2{#1/#2}\def\frac#1#2{#1/#2}%
  \def\vec#1{#1⃗}}%
\RequirePackage[official]{eurosym}

\tcbset{exercice corrige/.style={
  enhanced,breakable,before upper=\raggedright,
  colback=white,colframe=coulExo,
  boxrule=0.8pt,arc=2pt,
  left=6pt,right=6pt,top=8pt,bottom=5pt,
  before skip=9pt,after skip=4pt,
  coltitle=white,fonttitle=\bfseries\small,
  attach boxed title to top left={xshift=8pt,yshift=-\tcboxedtitleheight/2},
  boxed title style={colback=coulExo,arc=2pt,boxrule=0pt}}}
\newcounter{exo}
\newtcolorbox{exercice}[1][]{exercice corrige,
  phantom={\signet[exercice-\arabic{exo}]{2}{Exercice \arabic{exo}}},
  title={Exercice \theexo\ifx\relax#1\relax\else{} --- #1\fi}}
\newenvironment{exo}[1][\relax]{\refstepcounter{exo}\begin{exercice}[#1]}{\end{exercice}}

\geometry{top=1.8cm,bottom=1cm,left=1cm,right=1cm,headheight=14pt,footskip=0.6cm}
\renewcommand{\lycee@pied}{\thepage}

\renewcommand{\theexo}{\ifnum\value{exo}<10 0\fi\arabic{exo}}
\tcbset{exercice corrige/.style={
  enhanced,breakable,before upper=\raggedright,
  colback=white,colframe=coulExo,
  boxrule=0.9pt,arc=2pt,
  left=8pt,right=8pt,top=8pt,bottom=6pt,
  before skip=8pt,after skip=6pt,
  coltitle=white,fonttitle=\bfseries,
  attach boxed title to top left={xshift=8pt,yshift=-\tcboxedtitleheight/2},
  boxed title style={colback=coulExo,arc=2pt,boxrule=0pt}}}

\newcommand{\entetefiche}[2]{%
  \begin{center}
    \begin{tcolorbox}[enhanced,width=0.92\linewidth,colback=white,colframe=coulPart,
        boxrule=1.4pt,arc=3pt,halign=center,top=6pt,bottom=6pt]
      {\large\bfseries\color{coulPart} CORRIGÉ --- FICHE D'EXERCICES}\\[3pt]
      {\normalsize\bfseries #1}\\[2pt]
      {\normalsize #2}
    \end{tcolorbox}
  \end{center}}

\newcounter{partie}
\newcommand{\partie}[1]{%
  \stepcounter{partie}%
  \besoinplace{8\baselineskip}\vspace{4pt}%
  \begin{tcolorbox}[enhanced,colback=coulPart!8,colframe=coulPart,boxrule=0pt,
      leftrule=3pt,arc=0pt,left=6pt,right=4pt,top=3pt,bottom=3pt,
      before skip=6pt,after skip=2pt]
    \bfseries\color{coulPart}\Roman{partie} --- #1\signet{1}{\Roman{partie} --- #1}
  \end{tcolorbox}}
\newcommand{\partiesansnum}[1]{%
  \besoinplace{8\baselineskip}\vspace{4pt}%
  \begin{tcolorbox}[enhanced,colback=coulPart!8,colframe=coulPart,boxrule=0pt,
      leftrule=3pt,arc=0pt,left=6pt,right=4pt,top=3pt,bottom=3pt,
      before skip=6pt,after skip=2pt]
    \bfseries\color{coulPart}#1\signet{1}{#1}
  \end{tcolorbox}}

\setlist[questions,1]{label=\arabic*\textdegree),leftmargin=*,itemsep=3pt,topsep=3pt}
\setlist[questions,2]{label=\alph*),leftmargin=*,itemsep=2pt,topsep=2pt}
\setlist[questions,3]{label=\roman*),leftmargin=*,itemsep=1pt,topsep=1pt}
\newlist{expr}{enumerate}{1}
\setlist[expr]{label={\Alph*\ $=$},leftmargin=*,itemsep=2pt,topsep=3pt}

\newcommand{\res}[1]{{\color{coulRep}#1}}
\newcommand{\rem}[1]{\par\smallskip{\small\itshape\color{black!60}$\rightarrow$ #1}}
\newenvironment{calculs}{\par\smallskip\noindent$\displaystyle\begin{aligned}[t]}%
  {\end{aligned}$\par\smallskip}
\newcommand{\justif}[1]{\quad\text{\small\itshape\color{black!60}#1}}
\newcommand{\cas}[1]{\par\smallskip\noindent\textbf{\color{coulExo}#1}\nobreak}
\newcommand{\explique}[2]{\par\smallskip\noindent
  \begin{minipage}[t]{0.6\linewidth}\raggedright #1\end{minipage}\hfill
  \begin{minipage}[t]{0.37\linewidth}\raggedright\leavevmode
    {\small\itshape\color{black!60}#2}\end{minipage}\par}

\setlength{\parskip}{2pt}

\makeatother
\entete{Chapitre 6 --- Probabilités conditionnelles --- corrigé}{Première spécialité}

\tikzset{evt/.style={inner sep=2pt},
  proba/.style={font=\footnotesize,fill=white,inner sep=1pt},
  nom/.style={font=\footnotesize,inner sep=1.5pt,sloped},
  arbA/.style={},arbnA/.style={},arbAB/.style={},arbAnB/.style={},arbnAB/.style={},
  arbnAnB/.style={},chemin/.style={coulExo,line width=1.4pt}}
\newcommand{\arbcontraires}[2]{\def\arbnA{#1}\def\arbnB{#2}}
\newcommand{\arbfeuilles}[4]{\def\arbfAB{#1}\def\arbfAnB{#2}\def\arbfnAB{#3}\def\arbfnAnB{#4}}
\newcommand{\arbsansfeuilles}{\def\arbpasfeuilles{}}
\newcommand{\arbvaleurs}{\def\arbseules{}\arbsansfeuilles}
\newcommand{\arblab}[2]{\ifx\arbseules\undefined
  $#1\if\relax\detokenize{#2}\relax\else=#2\fi$\else$#2$\fi}
\newcommand{\arbfeuille}[2]{%
  \node[right,font=\footnotesize] at (#1.east) {$#2\ifx\arbtmp\empty\else=\arbtmp\fi$};}
\newcommand{\arbre}[8]{%
  \ifx\arbnA\undefined\def\arbNA{\overline{#1}}\else\let\arbNA\arbnA\fi
  \ifx\arbnB\undefined\def\arbNB{\overline{#2}}\else\let\arbNB\arbnB\fi
  \node (O) at (0,0) {};
  \node[evt] (A) at (3.4,1.2) {$#1$};
  \node[evt] (nA) at (3.4,-1.2) {$\arbNA$};
  \node[evt] (AB) at (7.2,1.9) {$#2$};
  \node[evt] (AnB) at (7.2,0.5) {$\arbNB$};
  \node[evt] (nAB) at (7.2,-0.5) {$#2$};
  \node[evt] (nAnB) at (7.2,-1.9) {$\arbNB$};
  \draw[arbA] (O.center) -- node[nom,above]{\arblab{P(#1)}{#3}} (A);
  \draw[arbnA] (O.center) -- node[nom,below]{\arblab{P(\arbNA)}{#4}} (nA);
  \draw[arbAB] (A) -- node[nom,above]{\arblab{P_{#1}(#2)}{#5}} (AB);
  \draw[arbAnB] (A) -- node[nom,below]{\arblab{P_{#1}(\arbNB)}{#6}} (AnB);
  \draw[arbnAB] (nA) -- node[nom,above]{\arblab{P_{\arbNA}(#2)}{#7}} (nAB);
  \draw[arbnAnB] (nA) -- node[nom,below]{\arblab{P_{\arbNA}(\arbNB)}{#8}} (nAnB);
  \ifx\arbpasfeuilles\undefined
    \ifx\arbfAB\undefined\arbfeuilles{}{}{}{}\fi
    \let\arbtmp\arbfAB\arbfeuille{AB}{P(#1\cap #2)}
    \let\arbtmp\arbfAnB\arbfeuille{AnB}{P(#1\cap\arbNB)}
    \let\arbtmp\arbfnAB\arbfeuille{nAB}{P(\arbNA\cap #2)}
    \let\arbtmp\arbfnAnB\arbfeuille{nAnB}{P(\arbNA\cap\arbNB)}
  \fi}
\newcommand{\feuille}[2]{\node[right,font=\footnotesize] at (#1.east) {#2};}
\newcommand{\arbremachines}[3]{%
  \node (O) at (0,0) {};
  \foreach \p/\q/\r [count=\i] in {#3}{%
    \pgfmathsetmacro{\y}{4.6-2.3*\i}
    \node[evt] (M\i) at (3.4,\y) {$#1_{\i}$};
    \ifnum\i=3 \tikzset{arbpos/.style={below}}\else\tikzset{arbpos/.style={above}}\fi
    \draw (O.center) -- node[nom,arbpos]{$P(#1_{\i})=\p$} (M\i);
    \node[evt] (D\i) at (7.2,\y+0.6) {$#2$};
    \node[evt] (nD\i) at (7.2,\y-0.6) {$\overline{#2}$};
    \draw (M\i) -- node[nom,above]{$P_{#1_{\i}}(#2)=\q$} (D\i);
    \draw (M\i) -- node[nom,below]{$P_{#1_{\i}}(\overline{#2})={\color{coulExo}\r}$} (nD\i);
    \node[right,font=\footnotesize] at (D\i.east) {$P(#1_{\i}\cap #2)$};
    \node[right,font=\footnotesize] at (nD\i.east) {$P(#1_{\i}\cap\overline{#2})$};}}
\newcommand{\arbretrois}[2]{%
  \node (O) at (0,0) {};
  \foreach \a/\ya/\na/\pa in {1/1.8/S/#1,0/-1.8/\overline{S}/#2}{%
    \node[evt] (N\a) at (2.2,\ya) {$\na$};
    \draw (O.center) -- node[proba]{$\pa$} (N\a);
    \foreach \b/\db/\nb/\pb in {1/0.9/S/#1,0/-0.9/\overline{S}/#2}{%
      \pgfmathsetmacro{\yb}{\ya+\db}
      \node[evt] (N\a\b) at (4.4,\yb) {$\nb$};
      \draw (N\a) -- node[proba]{$\pb$} (N\a\b);
      \foreach \c/\dc/\nc/\pc in {1/0.45/S/#1,0/-0.45/\overline{S}/#2}{%
        \pgfmathsetmacro{\yc}{\yb+\dc}
        \node[evt] (N\a\b\c) at (7,\yc) {$\nc$};
        \draw (N\a\b) -- node[proba]{$\pc$} (N\a\b\c);}}}}

\newcommand{\nv}[1]{{\color{coulPart}#1}}
\setlength{\parindent}{0pt}

\begin{document}

\entetefiche{Chapitre 6 --- Probabilités conditionnelles}{Première spécialité mathématiques}

\partie{Vocabulaire des probabilités (rappels de seconde)}

\begin{exo}[Décrire des événements]
\cas{1.}
\explique{$\Omega=\ens{1\,;2\,;\dots\,;24}$ : $24$ issues, toutes de même
probabilité.}{tirage au hasard : équiprobabilité}
\begin{calculs}
A &= \ens{4\,;8\,;12\,;16\,;20\,;24} &&\justif{multiples de $4$ jusqu'à $24$}\\
B &= \ens{6\,;12\,;18\,;24} &&\justif{multiples de $6$}\\
A\cap B &= \ens{12\,;24} &&\justif{dans $A$ \emph{et} dans $B$}\\
A\cup B &= \ens{4\,;6\,;8\,;12\,;16\,;18\,;20\,;24} &&\justif{dans $A$ \emph{ou} dans
$B$ (ou les deux)}
\end{calculs}
\cas{2.}
\begin{calculs}
P(A) &= \dfrac{6}{24} = \dfrac{1\times\barre{6}}{4\times\barre{6}} &&\justif{nombre
d'issues de $A$ sur $24$}\\
P(A) &= \res{\dfrac{1}{4}} &&\\
P(B) &= \dfrac{4}{24} = \dfrac{1\times\barre{4}}{6\times\barre{4}} &&\\
P(B) &= \res{\dfrac{1}{6}} &&\\
P(A\cap B) &= \dfrac{2}{24} = \res{\dfrac{1}{12}} &&\justif{$2$ issues}
\end{calculs}
\cas{3.}
\begin{calculs}
P(A\cup B) &= P(A)+P(B)-P(A\cap B) &&\justif{formule de la réunion}\\
P(A\cup B) &= \dfrac{1\fois{3}}{4\fois{3}}+\dfrac{1\fois{2}}{6\fois{2}}-\dfrac{1}{12}
  &&\justif{même dénominateur $12$}\\
P(A\cup B) &= \dfrac{3+2-1}{12} = \dfrac{4}{12} &&\\
P(A\cup B) &= \res{\dfrac{1}{3}} &&
\end{calculs}
\explique{En comptant : $A\cup B$ a $8$ issues, et $\dfrac{8}{24}=\dfrac{1}{3}$.}{on
retrouve le même résultat}
\rem{On retire $P(A\cap B)$ parce que $12$ et $24$ sont comptés deux fois dans
$P(A)+P(B)$.}
\cas{4.}
\explique{$\overline{A}$ : \og le nombre n'est pas un multiple de $4$\fg.}{contraire :
toutes les issues hors de $A$}
\begin{calculs}
P(\overline{A}) &= 1-P(A) &&\justif{probabilité du contraire}\\
P(\overline{A}) &= 1-\dfrac{1}{4} = \res{\dfrac{3}{4}} &&
\end{calculs}
\end{exo}

\begin{exo}[Avec les formules]
\cas{1.}
\explique{$S\cap M$ : \og l'élève pratique un sport en club \emph{et} joue d'un
instrument\fg.}{\og et\fg{} : intersection}
\explique{$S\cup M$ : \og l'élève pratique un sport en club \emph{ou} joue d'un instrument
(ou les deux)\fg, c'est-à-dire au moins l'une des deux activités.}{\og ou\fg{} :
réunion}
\cas{2.}
\begin{calculs}
P(S\cup M) &= P(S)+P(M)-P(S\cap M) &&\justif{formule de la réunion}\\
P(S\cup M) &= 0{,}45+0{,}3-0{,}15 &&\\
P(S\cup M) &= \res{0{,}6} &&
\end{calculs}
\cas{3.}
\explique{\og Ni sport en club ni musique\fg{} est le contraire de \og au moins l'une des
deux activités\fg, donc de $S\cup M$.}{$\overline{S\cup M}$}
\begin{calculs}
P(\overline{S\cup M}) &= 1-P(S\cup M) &&\justif{probabilité du contraire}\\
P(\overline{S\cup M}) &= 1-0{,}6 = \res{0{,}4} &&
\end{calculs}
\end{exo}

\begin{exo}[Retrouver une probabilité]
\cas{a)}
\explique{On part de la formule de la réunion et on isole $P(B)$.}{trois des quatre
nombres sont connus}
\begin{calculs}
P(A\cup B) &= P(A)+P(B)-P(A\cap B) &&\justif{formule de la réunion}\\
0{,}7 &= 0{,}4+P(B)-0{,}1 &&\justif{on remplace}\\
0{,}7 &= 0{,}3+P(B) &&\\
P(B) &= \res{0{,}4} &&\justif{on retire $0{,}3$}\\
P(\overline{B}) &= 1-0{,}4 = \res{0{,}6} &&\justif{probabilité du contraire}
\end{calculs}
\cas{b)}
\begin{calculs}
P(A) &= 1-P(\overline{A}) &&\justif{probabilité du contraire}\\
P(A) &= 1-0{,}35 = \res{0{,}65} &&\\
P(B) &= 1-0{,}5 = \res{0{,}5} &&\\
P(A\cap B) &= P(A)+P(B)-P(A\cup B) &&\justif{formule de la réunion, $P(A\cap B)$
isolée}\\
P(A\cap B) &= 0{,}65+0{,}5-0{,}8 &&\\
P(A\cap B) &= \res{0{,}35} &&
\end{calculs}
\cas{c)}
\explique{$A$ et $B$ sont incompatibles : $A\cap B=\varnothing$, donc
$P(A\cap B)=0$.}{ils ne peuvent pas se produire ensemble}
\begin{calculs}
P(A\cup B) &= P(A)+P(B) &&\justif{événements incompatibles}\\
P(A\cup B) &= 0{,}2+0{,}45 = \res{0{,}65} &&
\end{calculs}
\end{exo}

\begin{exo}[Deux activités dans une classe]
\cas{1.}\par\nobreak
\noindent\begin{minipage}[t]{0.5\linewidth}\vspace{0pt}\raggedright
On note $S$ : \og l'élève pratique un sport\fg{} et $I$ : \og il joue d'un
instrument\fg.
\par\smallskip On commence par l'intersection : $6$ élèves font les deux.
\par Sport seulement : $18-6=12$ ; instrument seulement : $10-6=4$.
\par Aucune des deux : $32-(12+6+4)=10$.
\end{minipage}\hfill
\begin{minipage}[t]{0.45\linewidth}\vspace{0pt}\centering
\begin{tikzpicture}[scale=1.2]
  \draw (-1.3,-1.05) rectangle (2.3,1.05);
  \node[font=\footnotesize] at (-1.1,0.85) {$\Omega$};
  \draw (0,0) circle (0.8);
  \draw (1,0) circle (0.8);
  \node[font=\footnotesize] at (-0.75,0.78) {$S$};
  \node[font=\footnotesize] at (1.75,0.78) {$I$};
  \node at (-0.35,0) {$12$};
  \node at (0.5,0) {$6$};
  \node at (1.35,0) {$4$};
  \node at (1.95,-0.8) {$10$};
\end{tikzpicture}
\end{minipage}
\cas{2.}
\explique{\og Au moins l'une des deux\fg{} : $S\cup I$, soit $12+6+4=22$ élèves.}{on
additionne les trois zones des cercles}
\begin{calculs}
P(S\cup I) &= \dfrac{22}{32} = \dfrac{11\times\barre{2}}{16\times\barre{2}}
  &&\justif{$22$ élèves sur $32$}\\
P(S\cup I) &= \res{\dfrac{11}{16}} &&
\end{calculs}
\cas{3.}
\begin{calculs}
P(\overline{S\cup I}) &= \dfrac{10}{32} = \res{\dfrac{5}{16}} &&\justif{la zone hors
des cercles}
\end{calculs}
\explique{On retrouve $1-\dfrac{11}{16}=\dfrac{5}{16}$.}{\og aucune\fg{} est le contraire
de \og au moins une\fg}
\cas{4.}
\explique{\og Exactement une\fg{} : sport seulement ou instrument seulement, soit
$12+4=16$ élèves.}{on exclut les $6$ élèves qui font les deux}
\begin{calculs}
P(\text{exactement une}) &= \dfrac{16}{32} = \res{\dfrac{1}{2}} &&
\end{calculs}
\end{exo}

\partie{Probabilités conditionnelles et tableaux}

\begin{exo}[Calculer dans un tableau]
\cas{1.}
\begin{calculs}
P(J) &= \dfrac{150}{300} = \res{\dfrac{1}{2}} &&\justif{total de la ligne
\og moins de $30$ ans\fg{} sur $300$}\\
P(V) &= \dfrac{90}{300} = \res{\dfrac{3}{10}} &&\justif{total de la colonne
\og Vélo\fg}\\
P(J\cap V) &= \dfrac{60}{300} = \res{\dfrac{1}{5}} &&\justif{case à la croisée de la
ligne et de la colonne}
\end{calculs}
\cas{2.}
\explique{On sait que le salarié a moins de $30$ ans : on cherche $P_{J}(V)$, et on ne
regarde que la ligne \og moins de $30$ ans\fg{} ($150$ salariés).}{la condition réduit
l'univers}
\begin{calculs}
P_{J}(V) &= \dfrac{60}{150} = \dfrac{2\times\barre{30}}{5\times\barre{30}}
  &&\justif{$60$ cyclistes parmi les $150$ moins de $30$ ans}\\
P_{J}(V) &= \res{\dfrac{2}{5}} &&
\end{calculs}
\cas{3.}
\explique{On sait qu'il vient à vélo : on cherche $P_{V}(J)$, dans la colonne
\og Vélo\fg{} ($90$ salariés).}{on change de condition}
\begin{calculs}
P_{V}(J) &= \dfrac{60}{90} = \dfrac{2\times\barre{30}}{3\times\barre{30}} &&\\
P_{V}(J) &= \res{\dfrac{2}{3}} &&
\end{calculs}
\rem{$P_{J}(V)\neq P_{V}(J)$ : l'ordre compte. Les deux probabilités ont le même
numérateur ($60$), mais pas le même dénominateur.}
\end{exo}

\begin{exo}[Compléter le tableau d'abord]
\cas{1.}
\explique{On place les données ($400$ ; $220$ ; $260$ ; $140$), puis on complète par
différences (en couleur) : hommes $400-220=180$ ; hommes reçus $260-140=120$ ;
femmes échouées $220-140=80$ ; hommes échoués $180-120=60$ ; échecs
$400-260=140$.}{chaque total est la somme de sa ligne ou de sa colonne}
\begin{center}
\renewcommand{\arraystretch}{1.3}
\begin{tabular}{|l|c|c|c|}
\hline
 & $R$ (reçu) & $\overline{R}$ (échec) & Total\\\hline
$F$ (femme) & $140$ & \nv{$80$} & $220$\\\hline
$\overline{F}$ (homme) & \nv{$120$} & \nv{$60$} & \nv{$180$}\\\hline
Total & $260$ & \nv{$140$} & $400$\\\hline
\end{tabular}
\end{center}
\cas{2.}
\begin{calculs}
P(\overline{F}\cap R) &= \dfrac{120}{400} = \res{\dfrac{3}{10}} &&\justif{hommes reçus
sur $400$}\\
P(F\cap\overline{R}) &= \dfrac{80}{400} = \res{\dfrac{1}{5}} &&\justif{femmes échouées}\\
P(\overline{R}) &= \dfrac{140}{400} = \res{\dfrac{7}{20}} &&\justif{total des échecs}
\end{calculs}
\cas{3.}
\begin{calculs}
P_{\overline{F}}(R) &= \dfrac{120}{180} = \dfrac{2\times\barre{60}}{3\times\barre{60}}
  &&\justif{parmi les $180$ hommes, $120$ reçus}\\
P_{\overline{F}}(R) &= \res{\dfrac{2}{3}} &&
\end{calculs}
\cas{4.}
\begin{calculs}
P_{\overline{R}}(F) &= \dfrac{80}{140} = \dfrac{4\times\barre{20}}{7\times\barre{20}}
  &&\justif{parmi les $140$ échecs, $80$ femmes}\\
P_{\overline{R}}(F) &= \res{\dfrac{4}{7}} &&
\end{calculs}
\end{exo}

\begin{exo}[Trois colonnes]
\cas{1.}
\begin{calculs}
P(\text{policier}) &= \dfrac{80}{200} = \res{\dfrac{2}{5}} &&\justif{total de la colonne
\og Policier\fg{} sur $200$}
\end{calculs}
\cas{2.}
\explique{On sait que l'élève est en première : on ne regarde que la ligne
\og Première\fg{} ($100$ élèves).}{probabilité conditionnelle}
\begin{calculs}
P_{\text{première}}(\text{policier}) &= \dfrac{50}{100} = \res{\dfrac{1}{2}} &&
\end{calculs}
\cas{3.}
\begin{calculs}
P_{\text{science-fiction}}(\text{seconde}) &= \dfrac{30}{60} = \res{\dfrac{1}{2}}
  &&\justif{colonne \og Science-fiction\fg{} : $30$ secondes sur $60$}
\end{calculs}
\cas{4.}
\explique{\og Parmi les élèves qui préfèrent les comédies\fg{} : colonne
\og Comédie\fg, $60$ élèves, dont $40$ en seconde ; $\dfrac{40}{60}=\dfrac{2}{3}$.}{on
calcule $P_{\text{comédie}}(\text{seconde})$}
\explique{L'affirmation est \res{vraie}.}{on conclut}
\end{exo}

\begin{exo}[Des pourcentages au tableau]
\cas{1.}
\explique{Site : $60\,\%$ de $2\,000$, soit $1\,200$ commandes ; application :
$800$. Renvoyées : $5\,\%$ de $1\,200$, soit $60$ ; $3\,\%$ de $800$, soit $24$.}{on
transforme chaque pourcentage en effectif}
\begin{center}
\renewcommand{\arraystretch}{1.3}
\begin{tabular}{|l|c|c|c|}
\hline
 & $R$ (renvoyée) & $\overline{R}$ & Total\\\hline
$W$ (site) & \nv{$60$} & \nv{$1\,140$} & \nv{$1\,200$}\\\hline
$\overline{W}$ (application) & \nv{$24$} & \nv{$776$} & \nv{$800$}\\\hline
Total & \nv{$84$} & \nv{$1\,916$} & $2\,000$\\\hline
\end{tabular}
\end{center}
\rem{$5\,\%$ porte sur les commandes passées sur le site, pas sur les $2\,000$ : c'est
une proportion \og parmi\fg, donc une probabilité conditionnelle, $P_{W}(R)=0{,}05$.}
\cas{2.}
\begin{calculs}
P(R) &= \dfrac{84}{2\,000} = \res{0{,}042} &&\justif{total de la colonne $R$}
\end{calculs}
\cas{3.}
\begin{calculs}
P_{R}(W) &= \dfrac{60}{84} = \dfrac{5\times\barre{12}}{7\times\barre{12}}
  &&\justif{parmi les $84$ renvoyées, $60$ du site}\\
P_{R}(W) &= \res{\dfrac{5}{7}} &&
\end{calculs}
\end{exo}

\partie{Formule des probabilités conditionnelles}

\begin{exo}[Traduire, puis calculer]
\cas{1.}
\begin{calculs}
P(I) &= 0{,}3 &&\justif{\og $30\,\%$ des élèves\fg{} : sur tout le lycée}\\
P(I\cap S) &= 0{,}12 &&\justif{\og internes \emph{et} inscrits\fg}\\
P_{\overline{I}}(S) &= 0{,}25 &&\justif{\og \emph{parmi} les élèves qui ne sont pas
internes\fg}
\end{calculs}
\cas{2.}
\begin{calculs}
P_{I}(S) &= \dfrac{P(I\cap S)}{P(I)} &&\justif{formule des probabilités
conditionnelles}\\
P_{I}(S) &= \dfrac{0{,}12}{0{,}3} &&\\
P_{I}(S) &= \res{0{,}4} &&
\end{calculs}
\explique{$40\,\%$ des internes sont inscrits à l'association sportive.}{on interprète :
\og parmi les internes\fg}
\cas{3.}
\explique{On cherche $P(\overline{I}\cap S)$, avec $P(\overline{I})=1-0{,}3=0{,}7$.}{\og
pas interne \emph{et} inscrit\fg}
\begin{calculs}
P(\overline{I}\cap S) &= P(\overline{I})\times P_{\overline{I}}(S) &&\justif{la formule,
écrite avec un produit}\\
P(\overline{I}\cap S) &= 0{,}7\times 0{,}25 &&\\
P(\overline{I}\cap S) &= \res{0{,}175} &&
\end{calculs}
\end{exo}

\begin{exo}[Appliquer la formule]
\cas{a)}
\begin{calculs}
P_{A}(B) &= \dfrac{P(A\cap B)}{P(A)} &&\justif{on divise par la probabilité de la
condition}\\
P_{A}(B) &= \dfrac{0{,}1}{0{,}5} = \res{0{,}2} &&\\
P_{B}(A) &= \dfrac{P(A\cap B)}{P(B)} &&\justif{la condition est maintenant $B$}\\
P_{B}(A) &= \dfrac{0{,}1}{0{,}4} = \res{0{,}25} &&
\end{calculs}
\cas{b)}
\begin{calculs}
P(A\cap B) &= P(A)\times P_{A}(B) &&\justif{la formule, écrite avec un produit}\\
P(A\cap B) &= 0{,}6\times 0{,}3 = \res{0{,}18} &&\\
P_{A}(\overline{B}) &= 1-P_{A}(B) &&\justif{contraire, sous la même condition $A$}\\
P_{A}(\overline{B}) &= 1-0{,}3 = \res{0{,}7} &&
\end{calculs}
\cas{c)}
\begin{calculs}
P_{B}(A) &= \dfrac{P(A\cap B)}{P(B)} &&\\
P_{B}(A) &= \dfrac{0{,}05}{0{,}25} = \res{0{,}2} &&\\
P_{B}(\overline{A}) &= 1-0{,}2 = \res{0{,}8} &&\justif{contraire, sous la condition $B$}
\end{calculs}
\cas{d)}
\begin{calculs}
P_{A}(B) &= \dfrac{P(A\cap B)}{P(A)} &&\\
P_{A}(B) &= \dfrac{1}{12}\div\dfrac{1}{3} = \dfrac{1}{12}\times 3 &&\justif{diviser par
$\frac{1}{3}$, c'est multiplier par $3$}\\
P_{A}(B) &= \dfrac{3}{12} = \dfrac{1\times\barre{3}}{4\times\barre{3}} &&\\
P_{A}(B) &= \res{\dfrac{1}{4}} &&
\end{calculs}
\cas{e)}
\begin{calculs}
P(A\cap B) &= P(A)\times P_{A}(B) &&\justif{la condition $A$ est connue}\\
P(A\cap B) &= 0{,}3\times 0{,}8 = \res{0{,}24} &&\\
P_{B}(A) &= \dfrac{P(A\cap B)}{P(B)} &&\justif{on écrit la formule avec la condition
$B$}\\
0{,}6 &= \dfrac{0{,}24}{P(B)} &&\justif{on remplace}\\
P(B) &= \dfrac{0{,}24}{0{,}6} = \res{0{,}4} &&\justif{on isole $P(B)$}
\end{calculs}
\end{exo}

\begin{exo}[\og Et\fg{} ou \og parmi\fg{} ?]
\cas{1.}
\begin{calculs}
P(F) &= 0{,}35 &&\justif{a) sur tous les inscrits}\\
P_{F}(L) &= 0{,}4 &&\justif{b) \og \emph{parmi} ceux qui empruntent des films\fg}\\
P(F\cap L) &= 0{,}14 &&\justif{c) \og des films \emph{et} des livres\fg}\\
P_{\overline{F}}(L) &= 0{,}5 &&\justif{d) \og un sur deux \emph{qui n'emprunte pas} de
films\fg}
\end{calculs}
\cas{2.}
\begin{calculs}
P(F\cap L) &= P(F)\times P_{F}(L) &&\justif{la formule, écrite avec un produit}\\
P(F\cap L) &= 0{,}35\times 0{,}4 &&\\
P(F\cap L) &= 0{,}14 &&
\end{calculs}
\explique{On retrouve bien la donnée c).}{les données sont cohérentes}
\cas{3.}
\explique{On cherche $P(\overline{F}\cap L)$, avec $P(\overline{F})=1-0{,}35=0{,}65$.}{\og
pas de films \emph{mais} des livres\fg}
\begin{calculs}
P(\overline{F}\cap L) &= P(\overline{F})\times P_{\overline{F}}(L) &&\\
P(\overline{F}\cap L) &= 0{,}65\times 0{,}5 = \res{0{,}325} &&
\end{calculs}
\end{exo}

\begin{exo}[Vrai ou faux]
\cas{1.}
\begin{calculs}
P_{A}(B) &= \dfrac{P(A\cap B)}{P(A)} &&\\
P_{A}(B) &= \dfrac{0{,}2}{0{,}4} = 0{,}5 &&
\end{calculs}
\explique{L'affirmation est \res{vraie}.}{}
\cas{2.}
\begin{calculs}
P_{B}(A) &= \dfrac{P(A\cap B)}{P(B)} &&\\
P_{B}(A) &= \dfrac{0{,}2}{0{,}5} = 0{,}4 &&
\end{calculs}
\explique{$0{,}4\neq 0{,}5$ : l'affirmation est \res{fausse}.}{même numérateur,
dénominateurs différents}
\cas{3.}
\begin{calculs}
P_{A}(\overline{B}) &= 1-P_{A}(B) &&\justif{contraire, sous la condition $A$}\\
P_{A}(\overline{B}) &= 1-0{,}5 = 0{,}5 &&
\end{calculs}
\explique{L'affirmation est \res{vraie}.}{}
\cas{4.}
\begin{calculs}
P(A\cup B) &= P(A)+P(B)-P(A\cap B) &&\justif{formule de la réunion}\\
P(A\cup B) &= 0{,}4+0{,}5-0{,}2 = 0{,}7 &&
\end{calculs}
\explique{$0{,}7\neq 0{,}9$ : l'affirmation est \res{fausse}.}{$0{,}9$ oublie de retirer
$P(A\cap B)$}
\end{exo}

\partie{Arbres pondérés}

\begin{exo}[Compléter un arbre]
\cas{1.}
\begin{calculs}
P(A) &= 0{,}4 &&\justif{première branche : probabilité de $A$}\\
P_{A}(B) &= 0{,}7 &&\justif{branche partant de $A$ : condition $A$}\\
P_{\overline{A}}(\overline{B}) &= 0{,}2 &&\justif{branche de $\overline{A}$ vers
$\overline{B}$}
\end{calculs}
\cas{2.}\par\nobreak
\noindent\begin{minipage}[t]{0.4\linewidth}\vspace{0pt}\raggedright
Règle du nœud : la somme des branches qui partent d'un même nœud vaut $1$.
\par\smallskip $P(\overline{A})=1-0{,}4=0{,}6$ ;
\par $P_{A}(\overline{B})=1-0{,}7=0{,}3$ ;
\par $P_{\overline{A}}(B)=1-0{,}2=0{,}8$.
\par\smallskip En couleur : complété.
\end{minipage}\hfill
\begin{minipage}[t]{0.58\linewidth}\vspace{0pt}\centering
\begin{tikzpicture}
  \arbfeuilles{}{\res{0{,}12}}{\res{0{,}48}}{}
  \arbre{A}{B}{0{,}4}{\nv{0{,}6}}{0{,}7}{\nv{0{,}3}}{\nv{0{,}8}}{0{,}2}
\end{tikzpicture}
\end{minipage}
\cas{3.}
\begin{calculs}
P(A\cap\overline{B}) &= P(A)\times P_{A}(\overline{B}) &&\justif{règle du chemin :
on multiplie le long des branches}\\
P(A\cap\overline{B}) &= 0{,}4\times 0{,}3 = \res{0{,}12} &&\\
P(\overline{A}\cap B) &= P(\overline{A})\times P_{\overline{A}}(B) &&\\
P(\overline{A}\cap B) &= 0{,}6\times 0{,}8 = \res{0{,}48} &&
\end{calculs}
\end{exo}

\begin{exo}[Construire un arbre]
\cas{1.}\par\nobreak
\noindent\begin{minipage}[t]{0.4\linewidth}\vspace{0pt}\raggedright
$P(C)=0{,}7$ (\og $70\,\%$ des trajets\fg) ;
\par $P_{C}(N)=0{,}2$ (\og parmi les trajets courts\fg) ;
\par $P_{\overline{C}}(N)=0{,}1$ (\og parmi les autres\fg).
\par\smallskip En couleur : les compléments à $1$ (règle du nœud).
\end{minipage}\hfill
\begin{minipage}[t]{0.58\linewidth}\vspace{0pt}\centering
\begin{tikzpicture}
  \arbfeuilles{\res{0{,}14}}{}{}{\res{0{,}27}}
  \arbre{C}{N}{0{,}7}{\nv{0{,}3}}{0{,}2}{\nv{0{,}8}}{0{,}1}{\nv{0{,}9}}
\end{tikzpicture}
\end{minipage}
\cas{2.}
\begin{calculs}
P(C\cap N) &= P(C)\times P_{C}(N) &&\justif{règle du chemin}\\
P(C\cap N) &= 0{,}7\times 0{,}2 = \res{0{,}14} &&
\end{calculs}
\explique{$14\,\%$ des trajets réservés sont des trajets courts annulés.}{on
interprète : \og et\fg, sur tous les trajets}
\cas{3.}
\begin{calculs}
P(\overline{C}\cap\overline{N}) &= P(\overline{C})\times P_{\overline{C}}(\overline{N})
  &&\justif{règle du chemin}\\
P(\overline{C}\cap\overline{N}) &= 0{,}3\times 0{,}9 = \res{0{,}27} &&
\end{calculs}
\explique{$27\,\%$ des trajets réservés ne sont pas courts et ne sont pas annulés.}{on
interprète}
\end{exo}

\begin{exo}[Une branche par la formule]
\cas{1.}
\begin{calculs}
P_{A}(B) &= \dfrac{P(A\cap B)}{P(A)} &&\justif{formule des probabilités
conditionnelles}\\
P_{A}(B) &= \dfrac{0{,}1}{0{,}25} = \res{0{,}4} &&
\end{calculs}
\cas{2.}\par\nobreak
\noindent\begin{minipage}[t]{0.4\linewidth}\vspace{0pt}\raggedright
$P(\overline{A})=1-0{,}25=0{,}75$ ;
\par $P_{A}(\overline{B})=1-0{,}4=0{,}6$ ;
\par $P_{\overline{A}}(\overline{B})=1-0{,}4=0{,}6$.
\par\smallskip En couleur : trouvé par le calcul ou par la règle du nœud.
\end{minipage}\hfill
\begin{minipage}[t]{0.58\linewidth}\vspace{0pt}\centering
\begin{tikzpicture}
  \arbfeuilles{0{,}1}{}{}{\res{0{,}45}}
  \arbre{A}{B}{0{,}25}{\nv{0{,}75}}{\nv{0{,}4}}{\nv{0{,}6}}{0{,}4}{\nv{0{,}6}}
\end{tikzpicture}
\end{minipage}
\cas{3.}
\begin{calculs}
P(\overline{A}\cap\overline{B}) &= P(\overline{A})\times P_{\overline{A}}(\overline{B})
  &&\justif{règle du chemin}\\
P(\overline{A}\cap\overline{B}) &= 0{,}75\times 0{,}6 = \res{0{,}45} &&
\end{calculs}
\end{exo}

\begin{exo}[Trouver l'erreur]
\cas{1.}
\explique{L'élève a écrit $0{,}18$ sur la branche de $C$ vers $A$. Or \og $18\,\%$ des
animaux sont des chiens adoptés\fg{} donne $P(C\cap A)=0{,}18$ : c'est une
intersection, sur tous les animaux, pas une probabilité \og parmi les chiens\fg.}{sur une
branche du second niveau : toujours une probabilité conditionnelle}
\cas{2.}
\begin{calculs}
P_{C}(A) &= \dfrac{P(C\cap A)}{P(C)} &&\justif{formule des probabilités
conditionnelles}\\
P_{C}(A) &= \dfrac{0{,}18}{0{,}6} = \res{0{,}3} &&
\end{calculs}
\begin{minipage}[t]{0.4\linewidth}\vspace{0pt}\raggedright
$P_{C}(\overline{A})=1-0{,}3=0{,}7$ ;
\par $P_{\overline{C}}(A)=0{,}4$ (\og parmi les chats\fg) et
$P_{\overline{C}}(\overline{A})=0{,}6$ : cette partie de l'arbre était juste.
\par\smallskip En couleur : corrigé.
\end{minipage}\hfill
\begin{minipage}[t]{0.58\linewidth}\vspace{0pt}\centering
\begin{tikzpicture}
  \arbfeuilles{0{,}18}{}{}{\res{0{,}24}}
  \arbre{C}{A}{0{,}6}{0{,}4}{\nv{0{,}3}}{\nv{0{,}7}}{0{,}4}{0{,}6}
\end{tikzpicture}
\end{minipage}
\cas{3.}
\explique{\og Un chat qui n'est pas adopté dans le mois\fg{} :
$\overline{C}\cap\overline{A}$.}{}
\begin{calculs}
P(\overline{C}\cap\overline{A}) &= P(\overline{C})\times P_{\overline{C}}(\overline{A})
  &&\justif{règle du chemin}\\
P(\overline{C}\cap\overline{A}) &= 0{,}4\times 0{,}6 = \res{0{,}24} &&
\end{calculs}
\end{exo}

\partie{Formule des probabilités totales}

\begin{exo}[Calculer $P(C)$, puis une probabilité inversée]
\cas{1.}\par\nobreak
\noindent\begin{minipage}[t]{0.4\linewidth}\vspace{0pt}\raggedright
$P(S)=0{,}4$ ;
\par $P_{S}(C)=0{,}75$ (\og parmi ceux qui viennent en semaine\fg) ;
\par $P_{\overline{S}}(C)=0{,}45$ (\og parmi les autres\fg).
\par\smallskip En couleur : les compléments à $1$ (règle du nœud).
\end{minipage}\hfill
\begin{minipage}[t]{0.58\linewidth}\vspace{0pt}\centering
\begin{tikzpicture}
  \arbfeuilles{\res{0{,}3}}{}{\res{0{,}27}}{}
  \arbre{S}{C}{0{,}4}{\nv{0{,}6}}{0{,}75}{\nv{0{,}25}}{0{,}45}{\nv{0{,}55}}
\end{tikzpicture}
\end{minipage}
\cas{2.}
\begin{calculs}
P(C) &= P(S\cap C)+P(\overline{S}\cap C) &&\justif{deux chemins mènent à $C$}\\
P(C) &= 0{,}4\times 0{,}75+0{,}6\times 0{,}45 &&\justif{règle du chemin sur chacun}\\
P(C) &= 0{,}3+0{,}27 &&\\
P(C) &= 0{,}57 &&
\end{calculs}
\cas{3.}
\explique{On cherche $P_{C}(S)$ : on \emph{sait} que l'adhérent prend un cours
collectif.}{probabilité \og inversée\fg{} : pas sur l'arbre}
\begin{calculs}
P_{C}(S) &= \dfrac{P(S\cap C)}{P(C)} &&\justif{formule des probabilités
conditionnelles}\\
P_{C}(S) &= \dfrac{0{,}3}{0{,}57} = \dfrac{30}{57} &&\justif{on multiplie en haut et en
bas par $100$}\\
P_{C}(S) &= \dfrac{10\times\barre{3}}{19\times\barre{3}} &&\\
P_{C}(S) &= \res{\dfrac{10}{19}} &&
\end{calculs}
\explique{$\dfrac{10}{19}\approx 0{,}53$ : environ $53\,\%$ des adhérents qui prennent un
cours collectif viennent en semaine.}{on interprète}
\end{exo}

\begin{exo}[Trois offres]
\cas{1.}\par\nobreak
\noindent\begin{minipage}[t]{0.4\linewidth}\vspace{0pt}\raggedright
Trois branches au premier niveau : $0{,}5$ ; $0{,}35$ et $1-0{,}5-0{,}35=0{,}15$ (règle
du nœud).
\par\smallskip Au second niveau, les taux de résiliation, puis leurs compléments à $1$, en
couleur.
\end{minipage}\hfill
\begin{minipage}[t]{0.58\linewidth}\vspace{0pt}\centering
\begin{tikzpicture}[y=0.8cm]
  \arbremachines{O}{R}{0{,}5/0{,}04/0{,}96,0{,}35/0{,}06/0{,}94,\nv{0{,}15}/0{,}1/0{,}9}
\end{tikzpicture}
\end{minipage}
\cas{2.}
\begin{calculs}
P(R) &= P(O_{1}\cap R)+P(O_{2}\cap R)+P(O_{3}\cap R) &&\justif{trois chemins mènent à
$R$}\\
P(R) &= 0{,}5\times 0{,}04+0{,}35\times 0{,}06+0{,}15\times 0{,}1 &&\justif{règle du
chemin}\\
P(R) &= 0{,}02+0{,}021+0{,}015 &&\\
P(R) &= \res{0{,}056} &&
\end{calculs}
\cas{3.}
\begin{calculs}
P_{R}(O_{3}) &= \dfrac{P(O_{3}\cap R)}{P(R)} &&\justif{probabilité inversée}\\
P_{R}(O_{3}) &= \dfrac{0{,}015}{0{,}056} = \res{\dfrac{15}{56}} &&\justif{on multiplie
par $1\,000$ ; $15=3\times 5$ et $56=2^{3}\times 7$ : irréductible}
\end{calculs}
\explique{$\dfrac{15}{56}\approx 0{,}27$ : environ $27\,\%$ des clients qui résilient ont
l'offre $O_{3}$, alors qu'elle ne concerne que $15\,\%$ des clients.}{on interprète}
\end{exo}

\begin{exo}[Une probabilité inconnue]
\cas{1.}\par\nobreak
\noindent\begin{minipage}[t]{0.4\linewidth}\vspace{0pt}\raggedright
$P(D)=0{,}4$ et $P_{D}(B)=0{,}6$ ;
\par $P_{\overline{D}}(B)=x$, la probabilité inconnue.
\par\smallskip En couleur : les compléments à $1$ (règle du nœud), dont $1-x$.
\end{minipage}\hfill
\begin{minipage}[t]{0.58\linewidth}\vspace{0pt}\centering
\begin{tikzpicture}
  \arbfeuilles{0{,}24}{}{0{,}6x}{}
  \arbre{D}{B}{0{,}4}{\nv{0{,}6}}{0{,}6}{\nv{0{,}4}}{x}{\nv{1-x}}
\end{tikzpicture}
\end{minipage}
\cas{2.}
\explique{$51\,\%$ des élèves prennent le bus : $P(B)=0{,}51$. On écrit $P(B)$ avec la
formule des probabilités totales.}{deux chemins mènent à $B$}
\begin{calculs}
P(B) &= P(D\cap B)+P(\overline{D}\cap B) &&\justif{probabilités totales}\\
P(B) &= 0{,}4\times 0{,}6+0{,}6\times x &&\justif{règle du chemin}\\
P(B) &= 0{,}24+0{,}6x &&
\end{calculs}
\explique{Comme $P(B)=0{,}51$ : $0{,}6x+0{,}24=0{,}51$.}{c'est l'égalité demandée}
\cas{3.}
\begin{calculs}
0{,}6x+0{,}24 &= 0{,}51 &&\justif{on recopie l'équation}\\
0{,}6x &= 0{,}27 &&\justif{on retire $0{,}24$}\\
x &= \dfrac{0{,}27}{0{,}6} &&\justif{on divise par $0{,}6$}\\
x &= \res{0{,}45} &&
\end{calculs}
\explique{$P_{\overline{D}}(B)=0{,}45$ : $45\,\%$ des élèves qui ne sont pas
demi-pensionnaires prennent le bus.}{on interprète $x$ : \og parmi\fg}
\end{exo}

\begin{exo}[Un test de dépistage]
\cas{1.}\par\nobreak
\noindent\begin{minipage}[t]{0.4\linewidth}\vspace{0pt}\raggedright
$P(M)=0{,}02$ ;
\par $P_{M}(T)=0{,}95$ (\og positif pour $95\,\%$ des malades\fg) ;
\par $P_{\overline{M}}(T)=0{,}1$.
\par\smallskip En couleur : les compléments à $1$.
\end{minipage}\hfill
\begin{minipage}[t]{0.58\linewidth}\vspace{0pt}\centering
\begin{tikzpicture}
  \arbfeuilles{\res{0{,}019}}{}{\res{0{,}098}}{\res{0{,}882}}
  \arbre{M}{T}{0{,}02}{\nv{0{,}98}}{0{,}95}{\nv{0{,}05}}{0{,}1}{\nv{0{,}9}}
\end{tikzpicture}
\end{minipage}
\cas{2.}
\begin{calculs}
P(T) &= P(M\cap T)+P(\overline{M}\cap T) &&\justif{deux chemins mènent à $T$}\\
P(T) &= 0{,}02\times 0{,}95+0{,}98\times 0{,}1 &&\justif{règle du chemin}\\
P(T) &= 0{,}019+0{,}098 &&\\
P(T) &= 0{,}117 &&
\end{calculs}
\cas{3.}
\begin{calculs}
P_{T}(M) &= \dfrac{P(M\cap T)}{P(T)} &&\justif{probabilité inversée}\\
P_{T}(M) &= \dfrac{0{,}019}{0{,}117} &&\\
P_{T}(M) &\approx \res{0{,}16} &&\justif{à la calculatrice, arrondi au centième}
\end{calculs}
\explique{Seulement $16\,\%$ environ des personnes testées positives sont malades.
Surprenant pour un test \og fiable à $95\,\%$\fg{} : la maladie est rare, et les
$10\,\%$ de faux positifs portent sur $98\,\%$ de la population.}{$0{,}098$ de faux
positifs contre $0{,}019$ de vrais}
\cas{4.}
\begin{calculs}
P(\overline{T}) &= 1-0{,}117 = 0{,}883 &&\justif{probabilité du contraire}\\
P(\overline{M}\cap\overline{T}) &= 0{,}98\times 0{,}9 = 0{,}882 &&\justif{règle du
chemin}\\
P_{\overline{T}}(\overline{M}) &= \dfrac{P(\overline{M}\cap\overline{T})}{P(\overline{T})}
  &&\justif{probabilité inversée}\\
P_{\overline{T}}(\overline{M}) &= \dfrac{0{,}882}{0{,}883} &&\\
P_{\overline{T}}(\overline{M}) &\approx \res{0{,}999} &&\justif{arrondi au millième}
\end{calculs}
\explique{Une personne testée négative n'est presque sûrement pas malade : un test
négatif est rassurant.}{on interprète}
\end{exo}

\partie{Événements indépendants}

\begin{exo}[Indépendants ou non ?]
\cas{1.}
\begin{calculs}
P(A) &= \dfrac{4}{8} = \dfrac{1}{2} &&\justif{$A=\ens{2\,;4\,;6\,;8}$}\\
P(B) &= \dfrac{4}{8} = \dfrac{1}{2} &&\justif{$B=\ens{1\,;2\,;3\,;4}$}\\
P(A\cap B) &= \dfrac{2}{8} = \dfrac{1}{4} &&\justif{$A\cap B=\ens{2\,;4}$}\\
P(A)\times P(B) &= \dfrac{1}{2}\times\dfrac{1}{2} = \dfrac{1}{4} &&\justif{on compare avec
$P(A\cap B)$}
\end{calculs}
\explique{$P(A\cap B)=P(A)\times P(B)$ : $A$ et $B$ sont \res{indépendants}.}{on
conclut}
\cas{2. a)}
\begin{calculs}
P(B) &= 0{,}3\times 0{,}6+0{,}7\times 0{,}6 &&\justif{probabilités totales}\\
P(B) &= 0{,}18+0{,}42 = 0{,}6 &&
\end{calculs}
\explique{$P_{A}(B)=0{,}6$ et $P(B)=0{,}6$ : $P_{A}(B)=P(B)$, donc $A$ et $B$ sont
\res{indépendants}.}{savoir que $A$ est réalisé ne change pas la probabilité de $B$}
\rem{On le voit directement sur l'arbre : $B$ a la même probabilité, $0{,}6$, après $A$
et après $\overline{A}$.}
\cas{2. b)}
\begin{calculs}
P(B) &= 0{,}5\times 0{,}2+0{,}5\times 0{,}4 &&\justif{probabilités totales}\\
P(B) &= 0{,}1+0{,}2 = 0{,}3 &&
\end{calculs}
\explique{$P_{A}(B)=0{,}2$ et $P(B)=0{,}3$ : $P_{A}(B)\neq P(B)$, donc $A$ et $B$ ne sont
\res{pas indépendants}.}{on lit $P_{A}(B)$ sur l'arbre}
\end{exo}

\begin{exo}[Retrouver $P(F)$]
\cas{1.}
\begin{calculs}
P(A\cap F) &= P(A)\times P(F) &&\justif{$A$ et $F$ sont indépendants}\\
P(A\cap F) &= 0{,}1p &&
\end{calculs}
\cas{2.}
\explique{\og Au moins l'un des deux défauts\fg{} : $A\cup F$, et $P(A\cup F)=0{,}28$.}{\og
au moins un\fg{} : réunion}
\begin{calculs}
P(A\cup F) &= P(A)+P(F)-P(A\cap F) &&\justif{formule de la réunion}\\
P(A\cup F) &= 0{,}1+p-0{,}1p &&\justif{on remplace}\\
P(A\cup F) &= 0{,}9p+0{,}1 &&\justif{$p-0{,}1p=0{,}9p$}
\end{calculs}
\explique{Comme $P(A\cup F)=0{,}28$ : $0{,}9p+0{,}1=0{,}28$.}{c'est l'égalité demandée}
\cas{3.}
\begin{calculs}
0{,}9p+0{,}1 &= 0{,}28 &&\justif{on recopie l'équation}\\
0{,}9p &= 0{,}18 &&\justif{on retire $0{,}1$}\\
p &= \dfrac{0{,}18}{0{,}9} &&\justif{on divise par $0{,}9$}\\
p &= 0{,}2 &&
\end{calculs}
\explique{$\res{P(F)=0{,}2}$.}{}
\end{exo}

\begin{exo}[Calculer avec l'indépendance]
\cas{1.}
\begin{calculs}
P(A\cap B) &= P(A)\times P(B) &&\justif{$A$ et $B$ sont indépendants}\\
P(A\cap B) &= 0{,}3\times 0{,}4 = \res{0{,}12} &&
\end{calculs}
\cas{2.}
\begin{calculs}
P(A\cup B) &= P(A)+P(B)-P(A\cap B) &&\justif{\og au moins un des deux jours\fg{} :
réunion}\\
P(A\cup B) &= 0{,}3+0{,}4-0{,}12 &&\\
P(A\cup B) &= \res{0{,}58} &&
\end{calculs}
\cas{3.}
\explique{\og Aucun des deux jours\fg{} est $\overline{A}\cap\overline{B}$ ;
$\overline{A}$ et $\overline{B}$ sont indépendants.}{contraires d'événements
indépendants}
\begin{calculs}
P(\overline{A}\cap\overline{B}) &= 0{,}7\times 0{,}6 &&\justif{$1-0{,}3$ et $1-0{,}4$}\\
P(\overline{A}\cap\overline{B}) &= \res{0{,}42} &&
\end{calculs}
\explique{On retrouve $1-0{,}58=0{,}42$ : \og aucun\fg{} est le contraire de \og au
moins un\fg.}{autre façon}
\end{exo}

\begin{exo}[Indépendance et second degré]
\cas{1.}
\begin{calculs}
P(A)\times P(B) &= x(1-x) &&\justif{on multiplie les deux probabilités}\\
P(A)\times P(B) &= x-x^{2} &&\justif{on développe}
\end{calculs}
\cas{2.}
\explique{Si $A$ et $B$ sont indépendants, alors $P(A\cap B)=P(A)\times P(B)$, soit
$0{,}21=x-x^{2}$.}{définition de l'indépendance}
\begin{calculs}
x-x^{2} &= 0{,}21 &&\justif{on recopie l'égalité}\\
-x^{2}+x-0{,}21 &= 0 &&\justif{on retire $0{,}21$}\\
x^{2}-x+0{,}21 &= 0 &&\justif{on multiplie par $-1$}
\end{calculs}
\cas{3.}
\explique{Équation du second degré : $a=1$, $b=-1$ et $c=0{,}21$.}{ch. 3}
\begin{calculs}
\Delta &= b^{2}-4ac &&\justif{discriminant}\\
\Delta &= (-1)^{2}-4\times 1\times 0{,}21 &&\\
\Delta &= 1-0{,}84 = 0{,}16 &&\justif{$\Delta>0$ : deux solutions ; $\sqrt{0{,}16}=0{,}4$}\\
x_{1} &= \dfrac{-b-\sqrt{\Delta}}{2a} = \dfrac{1-0{,}4}{2} &&\\
x_{1} &= 0{,}3 &&\\
x_{2} &= \dfrac{-b+\sqrt{\Delta}}{2a} = \dfrac{1+0{,}4}{2} &&\\
x_{2} &= 0{,}7 &&
\end{calculs}
\explique{Les deux solutions sont dans $\intff{0}{1}$ : $\res{P(A)=0{,}3}$ et
$\res{P(B)=0{,}7}$, ou $\res{P(A)=0{,}7}$ et $\res{P(B)=0{,}3}$.}{$P(B)=1-x$}
\rem{On vérifie : $0{,}3\times 0{,}7=0{,}21$.}
\end{exo}

\partie{Trois façons de tirer deux objets}

\begin{exo}[Avec remise, sans remise, simultané]
\explique{On note $R_{1}$ : \og la première boule est rouge\fg, $R_{2}$ : \og la seconde
boule est rouge\fg{} ; $V_{1}$ et $V_{2}$ de même.}{$\overline{R_{1}}=V_{1}$}
\cas{a)}\par\nobreak
\noindent\begin{minipage}[t]{0.4\linewidth}\vspace{0pt}\raggedright
Avec remise, l'urne ne change pas : $\frac{2}{4}=\frac{1}{2}$ pour chaque couleur, aux
deux niveaux.
\end{minipage}\hfill
\begin{minipage}[t]{0.58\linewidth}\vspace{0pt}\centering
\begin{tikzpicture}
  \arbcontraires{V_{1}}{V_{2}}\arbsansfeuilles
  \arbre{R_{1}}{R_{2}}{\frac{1}{2}}{\frac{1}{2}}{\frac{1}{2}}{\frac{1}{2}}{\frac{1}{2}}{\frac{1}{2}}
\end{tikzpicture}
\end{minipage}
\begin{calculs}
P(M) &= P(R_{1}\cap R_{2})+P(V_{1}\cap V_{2}) &&\justif{deux chemins : deux rouges ou
deux vertes}\\
P(M) &= \dfrac{1}{2}\times\dfrac{1}{2}+\dfrac{1}{2}\times\dfrac{1}{2} &&\justif{règle du
chemin}\\
P(M) &= \dfrac{1}{4}+\dfrac{1}{4} &&\\
P(M) &= \res{\dfrac{1}{2}} &&
\end{calculs}
\cas{b)}\par\nobreak
\noindent\begin{minipage}[t]{0.4\linewidth}\vspace{0pt}\raggedright
Sans remise, il reste $3$ boules : après une rouge, $1$ rouge et $2$ vertes ; après une
verte, $2$ rouges et $1$ verte.
\end{minipage}\hfill
\begin{minipage}[t]{0.58\linewidth}\vspace{0pt}\centering
\begin{tikzpicture}
  \arbcontraires{V_{1}}{V_{2}}\arbsansfeuilles
  \arbre{R_{1}}{R_{2}}{\frac{1}{2}}{\frac{1}{2}}{\frac{1}{3}}{\frac{2}{3}}{\frac{2}{3}}{\frac{1}{3}}
\end{tikzpicture}
\end{minipage}
\begin{calculs}
P(M) &= P(R_{1}\cap R_{2})+P(V_{1}\cap V_{2}) &&\justif{deux chemins}\\
P(M) &= \dfrac{1}{2}\times\dfrac{1}{3}+\dfrac{1}{2}\times\dfrac{1}{3} &&\justif{règle du
chemin}\\
P(M) &= \dfrac{1}{6}+\dfrac{1}{6} = \dfrac{2}{6} &&\\
P(M) &= \res{\dfrac{1}{3}} &&
\end{calculs}
\cas{c)}
\explique{On numérote les boules $r_{1}$, $r_{2}$ (rouges), $v_{1}$, $v_{2}$ (vertes) et
on liste les paires :}{sans ordre : $\ens{r_{1};r_{2}}=\ens{r_{2};r_{1}}$}
\begin{center}
$\ens{r_{1};r_{2}}$\quad $\ens{v_{1};v_{2}}$\quad $\ens{r_{1};v_{1}}$\quad
$\ens{r_{1};v_{2}}$\quad $\ens{r_{2};v_{1}}$\quad $\ens{r_{2};v_{2}}$
\end{center}
\begin{calculs}
P(M) &= \dfrac{2}{6} = \res{\dfrac{1}{3}} &&\justif{$2$ paires de même couleur sur $6$}
\end{calculs}
\explique{Même résultat que sans remise.}{simultané et sans remise se valent}
\end{exo}

\begin{exo}[Deux bonbons]
\cas{1.}\par\nobreak
\noindent\begin{minipage}[t]{0.4\linewidth}\vspace{0pt}\raggedright
Premier bonbon : $3$ au citron sur $5$.
\par\smallskip Léo mange le premier : il en reste $4$. Après un citron : $2$ citrons et
$2$ fraises ; après une fraise : $3$ citrons et $1$ fraise.
\end{minipage}\hfill
\begin{minipage}[t]{0.58\linewidth}\vspace{0pt}\centering
\begin{tikzpicture}
  \arbsansfeuilles
  \arbre{C_{1}}{C_{2}}{\frac{3}{5}}{\frac{2}{5}}{\frac{2}{4}}{\frac{2}{4}}{\frac{3}{4}}{\frac{1}{4}}
\end{tikzpicture}
\end{minipage}
\cas{2.}
\explique{\og Deux bonbons à la fraise\fg{} : $\overline{C_{1}}\cap\overline{C_{2}}$.}{le
chemin du bas}
\begin{calculs}
P(\overline{C_{1}}\cap\overline{C_{2}}) &= \dfrac{2}{5}\times\dfrac{1}{4} &&\justif{règle
du chemin}\\
P(\overline{C_{1}}\cap\overline{C_{2}}) &= \dfrac{2}{20} = \res{\dfrac{1}{10}} &&
\end{calculs}
\cas{3.}
\explique{\og Au moins un citron\fg{} est le contraire de \og deux fraises\fg.}{un seul
chemin à retirer}
\begin{calculs}
P(\text{au moins un citron}) &= 1-\dfrac{1}{10} = \res{\dfrac{9}{10}} &&\justif{probabilité
du contraire}
\end{calculs}
\cas{4.}
\explique{\og Un bonbon de chaque parfum\fg{} : citron puis fraise, ou fraise puis
citron.}{deux chemins}
\begin{calculs}
P(\text{un de chaque}) &= P(C_{1}\cap\overline{C_{2}})+P(\overline{C_{1}}\cap C_{2})
  &&\justif{deux chemins}\\
P(\text{un de chaque}) &= \dfrac{3}{5}\times\dfrac{2}{4}+\dfrac{2}{5}\times\dfrac{3}{4}
  &&\justif{règle du chemin}\\
P(\text{un de chaque}) &= \dfrac{6}{20}+\dfrac{6}{20} = \dfrac{12}{20} &&\\
P(\text{un de chaque}) &= \dfrac{3\times\barre{4}}{5\times\barre{4}} = \res{\dfrac{3}{5}}
  &&
\end{calculs}
\end{exo}

\begin{exo}[Un code de deux chiffres]
\cas{1.}
\explique{Avec répétition : $4$ choix pour le premier chiffre, $4$ pour le second, soit
$4\times 4=16$ codes, tous de même probabilité. Codes à deux chiffres pairs : $22$ ;
$24$ ; $42$ ; $44$.}{comme un tirage avec remise}
\begin{calculs}
P(E) &= \dfrac{4}{16} = \res{\dfrac{1}{4}} &&
\end{calculs}
\cas{2.}
\explique{Chiffres différents : $4$ choix, puis $3$, soit $4\times 3=12$ codes. Codes à
deux chiffres pairs : $24$ et $42$.}{comme un tirage sans remise}
\begin{calculs}
P(E) &= \dfrac{2}{12} = \res{\dfrac{1}{6}} &&
\end{calculs}
\cas{3.}
\explique{Les paires : $\ens{1;2}$, $\ens{1;3}$, $\ens{1;4}$, $\ens{2;3}$, $\ens{2;4}$,
$\ens{3;4}$ ; une seule est formée de deux chiffres pairs, $\ens{2;4}$.}{comme un tirage
simultané}
\begin{calculs}
P(E) &= \res{\dfrac{1}{6}} &&\justif{$1$ paire sur $6$}
\end{calculs}
\cas{4.}
\explique{Chaque paire donne exactement deux codes ($\ens{2;4}$ donne $24$ et $42$). Il
y a donc deux fois plus de codes que de paires, au total ($12$ contre $6$) comme dans $E$
($2$ contre $1$) : la proportion ne change pas.}{$\dfrac{2}{12}=\dfrac{1}{6}$ dans les deux cas}
\end{exo}

\partie{Épreuve de Bernoulli}

\begin{exo}[Trois tirs au but]
\cas{1.}\par\nobreak
\noindent\begin{minipage}[t]{0.4\linewidth}\vspace{0pt}\raggedright
Chaque tir est une épreuve de Bernoulli : succès $S$ (tir marqué) de probabilité
$0{,}6$, échec de probabilité $1-0{,}6=0{,}4$.
\par\smallskip Les tirs sont indépendants : les mêmes probabilités à chaque niveau.
\end{minipage}\hfill
\begin{minipage}[t]{0.58\linewidth}\vspace{0pt}\centering
\begin{tikzpicture}[y=0.8cm]
  \arbretrois{0{,}6}{0{,}4}
\end{tikzpicture}
\end{minipage}
\cas{2.}
\begin{calculs}
P(SSS) &= 0{,}6\times 0{,}6\times 0{,}6 &&\justif{règle du chemin}\\
P(SSS) &= \res{0{,}216} &&
\end{calculs}
\cas{3.}
\explique{Trois chemins ont exactement deux succès : $SS\overline{S}$,
$S\overline{S}S$ et $\overline{S}SS$.}{l'échec est au 1\textsuperscript{er},
2\textsuperscript{e} ou 3\textsuperscript{e} tir}
\begin{calculs}
P(\text{2 succès}) &= 3\times(0{,}6\times 0{,}6\times 0{,}4) &&\justif{trois chemins de
même probabilité}\\
P(\text{2 succès}) &= 3\times 0{,}144 = \res{0{,}432} &&
\end{calculs}
\cas{4.}
\explique{\og Au moins un tir marqué\fg{} est le contraire de \og aucun\fg{}
($\overline{S}\,\overline{S}\,\overline{S}$).}{un seul chemin à retirer}
\begin{calculs}
P(\text{au moins un}) &= 1-0{,}4\times 0{,}4\times 0{,}4 &&\justif{probabilité du
contraire}\\
P(\text{au moins un}) &= 1-0{,}064 = \res{0{,}936} &&
\end{calculs}
\end{exo}

\begin{exo}[Répondre au hasard]
\cas{1.}
\explique{Une réponse au hasard n'a que deux issues : juste (succès) ou fausse (échec).
Une seule réponse juste sur quatre, choisie au hasard : $\res{p=\dfrac{1}{4}}$.}{épreuve
de Bernoulli : deux issues}
\cas{2.}\par\nobreak
\noindent\begin{minipage}[t]{0.4\linewidth}\vspace{0pt}\raggedright
$S$ : \og la réponse est juste\fg ; succès $\frac{1}{4}$, échec $\frac{3}{4}$.
\par\smallskip Les trois réponses sont indépendantes : les mêmes probabilités à chaque
niveau.
\end{minipage}\hfill
\begin{minipage}[t]{0.58\linewidth}\vspace{0pt}\centering
\begin{tikzpicture}[y=0.8cm]
  \arbretrois{\frac{1}{4}}{\frac{3}{4}}
\end{tikzpicture}
\end{minipage}
\cas{3.}
\begin{calculs}
P(SSS) &= \dfrac{1}{4}\times\dfrac{1}{4}\times\dfrac{1}{4} = \res{\dfrac{1}{64}}
  &&\justif{règle du chemin}
\end{calculs}
\cas{4.}
\begin{calculs}
P(\overline{S}\,\overline{S}\,\overline{S}) &=
  \dfrac{3}{4}\times\dfrac{3}{4}\times\dfrac{3}{4} = \res{\dfrac{27}{64}} &&\justif{le
chemin du bas}
\end{calculs}
\cas{5.}
\explique{\og Au moins deux\fg{} : exactement deux réponses justes (trois chemins), ou
trois (un chemin).}{on additionne les chemins}
\begin{calculs}
P(\text{2 justes}) &= 3\times\dfrac{1}{4}\times\dfrac{1}{4}\times\dfrac{3}{4}
  &&\justif{la fausse au 1\textsuperscript{er}, 2\textsuperscript{e} ou
3\textsuperscript{e} rang}\\
P(\text{2 justes}) &= \dfrac{9}{64} &&\\
P(\text{au moins 2}) &= \dfrac{9}{64}+\dfrac{1}{64} = \dfrac{10}{64} &&\\
P(\text{au moins 2}) &= \dfrac{5\times\barre{2}}{32\times\barre{2}} = \res{\dfrac{5}{32}}
  &&
\end{calculs}
\end{exo}

\begin{exo}[Deux feux tricolores]
\cas{1.}\par\nobreak
\noindent\begin{minipage}[t]{0.4\linewidth}\vspace{0pt}\raggedright
$V_{1}$ : \og le premier feu est vert\fg, $V_{2}$ : \og le second feu est vert\fg.
\par\smallskip Les feux sont indépendants : $p$ et $1-p$ aux deux niveaux.
\end{minipage}\hfill
\begin{minipage}[t]{0.58\linewidth}\vspace{0pt}\centering
\begin{tikzpicture}
  \arbsansfeuilles
  \arbre{V_{1}}{V_{2}}{p}{1-p}{p}{1-p}{p}{1-p}
\end{tikzpicture}
\end{minipage}
\cas{2.}
\begin{calculs}
P(V_{1}\cap V_{2}) &= p\times p = \res{p^{2}} &&\justif{règle du chemin}
\end{calculs}
\cas{3.}
\begin{calculs}
p^{2} &= 0{,}36 &&\justif{on recopie l'égalité}\\
p &= \sqrt{0{,}36} = \res{0{,}6} &&\justif{$p$ est une probabilité : $p\geq 0$, on
écarte $-0{,}6$}
\end{calculs}
\cas{4.}
\explique{\og Un seul feu vert\fg{} : vert puis non vert, ou non vert puis vert.}{deux
chemins}
\begin{calculs}
P(\text{un seul vert}) &= p(1-p)+(1-p)p &&\justif{règle du chemin}\\
P(\text{un seul vert}) &= 2\times 0{,}6\times 0{,}4 &&\justif{avec $p=0{,}6$}\\
P(\text{un seul vert}) &= \res{0{,}48} &&
\end{calculs}
\end{exo}

\partie{Variable aléatoire}

\begin{exo}[Un jeu de dé]
\cas{1.}
\explique{$5$ ou $6$ : gain $4$ ; $4$ : gain $1$ ; $1$ ; $2$ ou $3$ : perte de $3$\,€,
soit un gain de $-3$. $X$ prend les valeurs $\res{-3}$ ; $\res{1}$ et $\res{4}$.}{une
perte est un gain négatif}
\cas{2.}
\begin{calculs}
P(X=4) &= \dfrac{2}{6} = \res{\dfrac{1}{3}} &&\justif{faces $5$ et $6$}\\
P(X=-3) &= \dfrac{3}{6} = \res{\dfrac{1}{2}} &&\justif{faces $1$ ; $2$ et $3$}
\end{calculs}
\cas{3.}
\explique{$X\geq 1$ : $X=1$ ou $X=4$, soit les faces $4$ ; $5$ et $6$.}{les valeurs de $X$
supérieures ou égales à $1$}
\begin{calculs}
P(X\geq 1) &= P(X=1)+P(X=4) &&\\
P(X\geq 1) &= \dfrac{1}{6}+\dfrac{2}{6} = \dfrac{3}{6} &&\\
P(X\geq 1) &= \res{\dfrac{1}{2}} &&
\end{calculs}
\cas{4.}
\explique{$\ens{X\leq 1}$ : \og le joueur gagne au plus $1$\,€\fg, c'est-à-dire qu'il
obtient $1$ ; $2$ ; $3$ ou $4$.}{$X=-3$ ou $X=1$}
\begin{calculs}
P(X\leq 1) &= \dfrac{4}{6} = \res{\dfrac{2}{3}} &&\justif{$4$ faces sur $6$}
\end{calculs}
\end{exo}

\begin{exo}[Un cube peint]
\cas{1.}
\explique{Trois faces peintes : les \res{$8$} coins du grand cube.}{un par sommet}
\explique{Deux faces peintes : le cube du milieu de chaque arête, soit \res{$12$}
cubes.}{un cube a $12$ arêtes}
\explique{Une seule face peinte : le cube du centre de chaque face, soit
\res{$6$} cubes.}{un cube a $6$ faces}
\explique{Aucune face peinte : le cube du centre, caché, soit \res{$1$} cube.}{on vérifie :
$8+12+6+1=27$}
\cas{2.}
\explique{$X$ prend les valeurs $\res{0}$ ; $\res{1}$ ; $\res{2}$ et $\res{3}$.}{}
\cas{3.}
\begin{calculs}
P(X=2) &= \dfrac{12}{27} = \dfrac{4\times\barre{3}}{9\times\barre{3}} &&\justif{$12$ cubes
sur $27$, tirés au hasard}\\
P(X=2) &= \res{\dfrac{4}{9}} &&\\
P(X\geq 1) &= 1-P(X=0) &&\justif{\og au moins une\fg{} : contraire de \og aucune\fg}\\
P(X\geq 1) &= 1-\dfrac{1}{27} = \res{\dfrac{26}{27}} &&
\end{calculs}
\cas{4.}
\explique{$\ens{X\geq 2}$ : \og le petit cube a au moins deux faces peintes\fg{} (un coin
ou un milieu d'arête).}{$X=2$ ou $X=3$}
\begin{calculs}
P(X\geq 2) &= \dfrac{12+8}{27} = \res{\dfrac{20}{27}} &&
\end{calculs}
\end{exo}

\begin{exo}[La somme de deux jetons]
\cas{1.}
\begin{center}
\renewcommand{\arraystretch}{1.3}
\begin{tabular}{|c|c|c|c|c|c|c|}
\hline
Paire & $\ens{1;2}$ & $\ens{1;3}$ & $\ens{1;4}$ & $\ens{2;3}$ & $\ens{2;4}$ & $\ens{3;4}$\\\hline
$X$ & $3$ & $4$ & $5$ & $5$ & $6$ & $7$\\\hline
\end{tabular}
\end{center}
\rem{Tirage simultané : $\ens{1;2}$ et $\ens{2;1}$ sont la même paire.}
\cas{2.}
\explique{$X$ prend les valeurs $\res{3}$ ; $\res{4}$ ; $\res{5}$ ; $\res{6}$ et
$\res{7}$.}{}
\cas{3.}
\begin{calculs}
P(X=5) &= \dfrac{2}{6} = \res{\dfrac{1}{3}} &&\justif{$\ens{1;4}$ et $\ens{2;3}$}\\
P(X\geq 6) &= \dfrac{2}{6} = \res{\dfrac{1}{3}} &&\justif{$\ens{2;4}$ et $\ens{3;4}$}
\end{calculs}
\cas{4.}
\begin{calculs}
P(\text{$X$ pair}) &= P(X=4)+P(X=6) &&\justif{$4$ et $6$ sont les valeurs paires}\\
P(\text{$X$ pair}) &= \dfrac{1}{6}+\dfrac{1}{6} = \res{\dfrac{1}{3}} &&
\end{calculs}
\end{exo}

\partie{Loi de probabilité}

\begin{exo}[Un jeu avec mise]
\cas{1.}
\explique{On peut tirer $0$ ; $1$ ou $2$ jetons gagnants. On reçoit alors $0$ ; $5$ ou
$10$\,€, et on retire les $3$\,€ payés : $X$ prend les valeurs $\res{-3}$ ; $\res{2}$ et
$\res{7}$.}{gain algébrique : reçu moins mise}
\cas{2.}
\explique{Avec remise : à chaque tirage, $P(G)=\dfrac{2}{10}=0{,}2$ (jeton gagnant) et
$0{,}8$ sinon, aux deux niveaux.}{$G_{1}$, $G_{2}$ : premier, second jeton gagnant}
\begin{calculs}
P(X=-3) &= 0{,}8\times 0{,}8 = 0{,}64 &&\justif{aucun gagnant : un chemin}\\
P(X=2) &= 0{,}2\times 0{,}8+0{,}8\times 0{,}2 = 0{,}32 &&\justif{un gagnant : deux
chemins}\\
P(X=7) &= 0{,}2\times 0{,}2 = 0{,}04 &&\justif{deux gagnants : un chemin}
\end{calculs}
\begin{center}
\renewcommand{\arraystretch}{1.4}
\begin{tabular}{|c|c|c|c|}
\hline
$x_{i}$ & $-3$ & $2$ & $7$\\\hline
$P(X=x_{i})$ & $\res{0{,}64}$ & $\res{0{,}32}$ & $\res{0{,}04}$\\\hline
\end{tabular}
\end{center}
\rem{On vérifie que la somme vaut $1$ : $0{,}64+0{,}32+0{,}04=1$.}
\end{exo}

\begin{exo}[Compléter une loi]
\cas{1.}
\explique{La somme des probabilités d'une loi vaut $1$.}{}
\begin{calculs}
0{,}15+a+0{,}4+0{,}1 &= 1 &&\justif{somme des probabilités}\\
a+0{,}65 &= 1 &&\\
a &= \res{0{,}35} &&\\
P(Y\geq 0) &= P(Y=0)+P(Y=1)+P(Y=4) &&\justif{les valeurs $0$ ; $1$ et $4$}\\
P(Y\geq 0) &= 0{,}35+0{,}4+0{,}1 = \res{0{,}85} &&\\
P(Y>0) &= P(Y=1)+P(Y=4) &&\justif{$0$ n'est pas strictement positif}\\
P(Y>0) &= 0{,}4+0{,}1 = \res{0{,}5} &&
\end{calculs}
\cas{2.}
\begin{calculs}
0{,}4+a+2a &= 1 &&\justif{somme des probabilités}\\
3a &= 0{,}6 &&\justif{on retire $0{,}4$}\\
a &= \res{0{,}2} &&\justif{on divise par $3$}
\end{calculs}
\rem{On vérifie : $0{,}4+0{,}2+0{,}4=1$.}
\end{exo}

\begin{exo}[Le prix d'une entrée]
\cas{1.}\par\nobreak
\noindent\begin{minipage}[t]{0.4\linewidth}\vspace{0pt}\raggedright
$P(A)=0{,}7$ ; $P_{A}(E)=0{,}3$ (\og parmi eux\fg) ; $P_{\overline{A}}(E)=0{,}5$.
\par\smallskip En couleur : les compléments à $1$.
\end{minipage}\hfill
\begin{minipage}[t]{0.58\linewidth}\vspace{0pt}\centering
\begin{tikzpicture}
  \arbfeuilles{\res{0{,}21}}{\res{0{,}49}}{0{,}15}{0{,}15}
  \arbre{A}{E}{0{,}7}{\nv{0{,}3}}{0{,}3}{\nv{0{,}7}}{0{,}5}{\nv{0{,}5}}
\end{tikzpicture}
\end{minipage}
\cas{2.}
\explique{Sans pass ($\overline{A}$) : $30$\,€, enfant ou adulte. Adulte avec pass
($A\cap\overline{E}$) : $10$\,€. Enfant avec pass ($A\cap E$) : $0$\,€. $X$ prend les
valeurs $\res{0}$ ; $\res{10}$ et $\res{30}$.}{une valeur par cas de l'énoncé}
\cas{3.}
\begin{calculs}
P(X=0) &= P(A\cap E) &&\justif{enfant avec pass}\\
P(X=0) &= 0{,}7\times 0{,}3 = 0{,}21 &&\justif{règle du chemin}\\
P(X=10) &= P(A\cap\overline{E}) &&\justif{adulte avec pass}\\
P(X=10) &= 0{,}7\times 0{,}7 = 0{,}49 &&\\
P(X=30) &= P(\overline{A}) = 0{,}3 &&\justif{les deux chemins de $\overline{A}$ :
$0{,}15+0{,}15$}
\end{calculs}
\begin{center}
\renewcommand{\arraystretch}{1.4}
\begin{tabular}{|c|c|c|c|}
\hline
$x_{i}$ & $0$ & $10$ & $30$\\\hline
$P(X=x_{i})$ & $\res{0{,}21}$ & $\res{0{,}49}$ & $\res{0{,}3}$\\\hline
\end{tabular}
\end{center}
\rem{On vérifie : $0{,}21+0{,}49+0{,}3=1$.}
\end{exo}

\begin{exo}[Un dé truqué]
\cas{1.}
\explique{Les faces $1$ à $6$ ont pour probabilités $q$, $2q$, $3q$, $4q$, $5q$ et
$6q$ ; leur somme vaut $1$.}{somme des probabilités d'une loi}
\begin{calculs}
q+2q+3q+4q+5q+6q &= 1 &&\justif{somme des probabilités}\\
21q &= 1 &&\justif{$1+2+3+4+5+6=21$}
\end{calculs}
\cas{2.}
\explique{$q=\dfrac{1}{21}$, et la face $k$ a pour probabilité $\dfrac{k}{21}$.}{on divise
par $21$}
\begin{center}
\renewcommand{\arraystretch}{2.4}
\begin{tabular}{|c|c|c|c|c|c|c|}
\hline
$x_{i}$ & $1$ & $2$ & $3$ & $4$ & $5$ & $6$\\\hline
$P(X=x_{i})$ & $\res{\dfrac{1}{21}}$ & $\res{\dfrac{2}{21}}$ & $\res{\dfrac{3}{21}}$ &
$\res{\dfrac{4}{21}}$ & $\res{\dfrac{5}{21}}$ & $\res{\dfrac{6}{21}}$\\\hline
\end{tabular}
\end{center}
\cas{3.}
\begin{calculs}
P(\text{pair}) &= P(X=2)+P(X=4)+P(X=6) &&\justif{faces paires}\\
P(\text{pair}) &= \dfrac{2+4+6}{21} = \dfrac{12}{21} &&\\
P(\text{pair}) &= \dfrac{4\times\barre{3}}{7\times\barre{3}} = \res{\dfrac{4}{7}} &&
\end{calculs}
\end{exo}

\partie{Espérance, variance, écart-type}

\begin{exo}[Le jeu est-il favorable ?]
\cas{1.}
\begin{calculs}
E(X) &= 0{,}64\times(-3)+0{,}32\times 2+0{,}04\times 7 &&\justif{somme des
$p_{i}\,x_{i}$}\\
E(X) &= -1{,}92+0{,}64+0{,}28 &&\\
E(X) &= \res{-1} &&
\end{calculs}
\cas{2.}
\explique{Sur un grand nombre de parties, le joueur perd en moyenne $1$\,€ par partie :
le jeu lui est \res{défavorable}.}{$E(X)<0$}
\cas{3.}
\explique{Les gains deviennent $-m$ ; $5-m$ ; $10-m$, avec les mêmes
probabilités.}{on retire $m$ au lieu de $3$}
\begin{calculs}
E(X) &= 0{,}64\times(-m)+0{,}32\times(5-m)+0{,}04\times(10-m) &&\justif{même formule}\\
E(X) &= -0{,}64m+1{,}6-0{,}32m+0{,}4-0{,}04m &&\justif{on développe}\\
E(X) &= 2-m &&\justif{on réduit}
\end{calculs}
\explique{Le jeu est équitable quand $2-m=0$, soit pour une mise de
\res{$2$\,€}.}{$E(X)=0$}
\end{exo}

\begin{exo}[Variance et écart-type]
\cas{Résolution}
\begin{calculs}
V(X) &= 0{,}64\times(-3+1)^{2}+0{,}32\times(2+1)^{2}+0{,}04\times(7+1)^{2}
  &&\justif{écarts à $E(X)=-1$, au carré}\\
V(X) &= 0{,}64\times 4+0{,}32\times 9+0{,}04\times 64 &&\\
V(X) &= 2{,}56+2{,}88+2{,}56 &&\\
V(X) &= \res{8} &&\\
\sigma(X) &= \sqrt{8} &&\justif{racine carrée de la variance}\\
\sigma(X) &\approx \res{2{,}83} &&\justif{arrondi au centième}
\end{calculs}
\rem{Erreur fréquente : écrire $(-3-1)^{2}$. L'écart est $x_{i}-E(X)$, ici
$-3-(-1)=-2$.}
\end{exo}

\begin{exo}[Deux jeux à comparer]
\cas{1.}
\begin{calculs}
E(X_{A}) &= 0{,}5\times 2+0{,}5\times(-2) &&\justif{somme des $p_{i}\,x_{i}$}\\
E(X_{A}) &= 1-1 = \res{0} &&\\
E(X_{B}) &= 0{,}1\times 18+0{,}9\times(-2) &&\justif{on perd avec la probabilité
$1-0{,}1=0{,}9$}\\
E(X_{B}) &= 1{,}8-1{,}8 = \res{0} &&
\end{calculs}
\explique{Les deux jeux sont équitables : en moyenne, on ne gagne ni ne perd.}{$E=0$}
\cas{2.}
\begin{calculs}
V(X_{A}) &= 0{,}5\times(2-0)^{2}+0{,}5\times(-2-0)^{2} &&\justif{écarts à $E(X_{A})=0$}\\
V(X_{A}) &= 2+2 = 4 &&\\
\sigma(X_{A}) &= \sqrt{4} = \res{2} &&\\
V(X_{B}) &= 0{,}1\times(18-0)^{2}+0{,}9\times(-2-0)^{2} &&\justif{écarts à
$E(X_{B})=0$}\\
V(X_{B}) &= 32{,}4+3{,}6 = 36 &&\\
\sigma(X_{B}) &= \sqrt{36} = \res{6} &&
\end{calculs}
\cas{3.}
\explique{Même espérance, mais $\sigma(X_{B})>\sigma(X_{A})$ : les gains du jeu B sont
plus dispersés (on perd souvent, on gagne rarement beaucoup). Le jeu \res{B} est le plus
risqué.}{l'écart-type mesure le risque}
\end{exo}

\begin{exo}[La tombola de la fête du lycée]
\cas{1.}
\explique{On retire les $2$\,€ du billet à chaque lot : $100-2=98$ ; $20-2=18$ ;
$5-2=3$ ; sans lot, $-2$. $X$ prend les valeurs $\res{-2}$ ; $\res{3}$ ; $\res{18}$ et
$\res{98}$.}{gain algébrique : lot moins prix}
\cas{2.}
\explique{Billets gagnants : $1+4+10=15$ ; perdants : $200-15=185$.}{chaque billet a la
probabilité $\frac{1}{200}$}
\begin{center}
\renewcommand{\arraystretch}{1.4}
\begin{tabular}{|c|c|c|c|c|}
\hline
$x_{i}$ & $-2$ & $3$ & $18$ & $98$\\\hline
$P(X=x_{i})$ & $\res{0{,}925}$ & $\res{0{,}05}$ & $\res{0{,}02}$ & $\res{0{,}005}$\\\hline
\end{tabular}
\end{center}
\rem{$\frac{185}{200}=0{,}925$ ; $\frac{10}{200}=0{,}05$ ; $\frac{4}{200}=0{,}02$ ;
$\frac{1}{200}=0{,}005$ ; la somme vaut $1$.}
\cas{3.}
\begin{calculs}
E(X) &= 0{,}925\times(-2)+0{,}05\times 3+0{,}02\times 18+0{,}005\times 98
  &&\justif{somme des $p_{i}\,x_{i}$}\\
E(X) &= -1{,}85+0{,}15+0{,}36+0{,}49 &&\\
E(X) &= \res{-0{,}85} &&
\end{calculs}
\explique{En moyenne, un participant perd $0{,}85$\,€ par billet : le jeu est
défavorable, et l'association gagne en moyenne $0{,}85$\,€ par billet vendu.}{on
interprète}
\cas{4.}
\explique{On vend $N$ billets. Les lots valent en tout $100+4\times 20+10\times
5=230$\,€ ; les participants paient $2N$ euros.}{l'espérance de gain est ce que reçoit en
moyenne un billet, moins son prix}
\begin{calculs}
E(X) &= \dfrac{230}{N}-2 &&\justif{lots partagés entre $N$ billets, moins le prix}\\
\dfrac{230}{N}-2 &= 0 &&\justif{jeu équitable}\\
\dfrac{230}{N} &= 2 &&\\
N &= \dfrac{230}{2} = \res{115} &&
\end{calculs}
\explique{Avec $115$ billets vendus, la recette ($230$\,€) paie exactement les lots : le
jeu est équitable.}{on vérifie : $115\times 2=230$}
\rem{Non exigé : la même formule donne, pour $N=200$, $\dfrac{230}{200}-2=-0{,}85$, le
résultat de la question 3.}
\end{exo}

\begin{exo}[Une valeur inconnue]
\cas{1.}
\begin{calculs}
E(X) &= 0{,}3\times(-2)+0{,}5\times a+0{,}2\times 3 &&\justif{somme des $p_{i}\,x_{i}$}\\
E(X) &= -0{,}6+0{,}5a+0{,}6 &&\\
E(X) &= \res{0{,}5a} &&
\end{calculs}
\cas{2.}
\begin{calculs}
0{,}5a &= 2 &&\justif{$E(X)=2$}\\
a &= \dfrac{2}{0{,}5} = \res{4} &&\justif{on divise par $0{,}5$}
\end{calculs}
\cas{3.}
\begin{calculs}
V(X) &= 0{,}3\times(-2-2)^{2}+0{,}5\times(4-2)^{2}+0{,}2\times(3-2)^{2}
  &&\justif{écarts à $E(X)=2$, au carré}\\
V(X) &= 0{,}3\times 16+0{,}5\times 4+0{,}2\times 1 &&\\
V(X) &= 4{,}8+2+0{,}2 &&\\
V(X) &= \res{7} &&\\
\sigma(X) &= \sqrt{7} \approx \res{2{,}65} &&\justif{arrondi au centième}
\end{calculs}
\end{exo}

\partiesansnum{Exercices type bac}

\begin{exo}[Automatismes (QCM)]
\cas{1.}
\explique{\og Strictement supérieur à $4$\fg{} : $5$ ou $6$, soit $\dfrac{2}{6}=\dfrac{1}{3}$.
Réponse \res{b)}.}{$4$ n'est pas compris}
\cas{2.}
\explique{Parmi les $12$ filles, $3$ internes : $\dfrac{3}{12}=\dfrac{1}{4}$. Réponse
\res{c)}.}{on choisit une fille : l'univers, ce sont les $12$ filles}
\cas{3.}
\explique{$P(\overline{A}\cap B)=P(\overline{A})\times P_{\overline{A}}(B)=0{,}7\times
0{,}5=0{,}35$. Réponse \res{b)}.}{règle du chemin ; $0{,}06$ est
$P(A\cap B)$}
\cas{4.}
\explique{$P(A\cap B)=0{,}5\times 0{,}4=0{,}2$, puis
$P(A\cup B)=0{,}5+0{,}4-0{,}2=0{,}7$. Réponse \res{c)}.}{indépendance, puis formule de la
réunion}
\cas{5.}
\explique{$E(X)=0{,}5\times(-1)+0{,}3\times 2+0{,}2\times 5=-0{,}5+0{,}6+1=1{,}1$.
Réponse \res{a)}.}{somme des $p_{i}\,x_{i}$}
\end{exo}

\begin{exo}[Location de kayaks]
\cas{1.}\par\nobreak
\noindent\begin{minipage}[t]{0.4\linewidth}\vspace{0pt}\raggedright
$P(S)=0{,}7$, donc $P(\overline{S})=0{,}3$ ;
\par $P_{S}(L)=0{,}2$ (\og parmi les locations en semaine\fg), donc
$P_{S}(\overline{L})=0{,}8$.
\par\smallskip Les branches issues de $\overline{S}$ restent vides : la question 4 les
calcule.
\end{minipage}\hfill
\begin{minipage}[t]{0.58\linewidth}\vspace{0pt}\centering
\begin{tikzpicture}
  \arbsansfeuilles
  \arbre{S}{L}{\nv{0{,}7}}{\nv{0{,}3}}{\nv{0{,}2}}{\nv{0{,}8}}{}{}
\end{tikzpicture}
\end{minipage}
\cas{2.}
\explique{\og $23\,\%$ de l'ensemble des locations\fg{} : $\res{P(L)=0{,}23}$.}{sur toutes
les locations : pas une probabilité conditionnelle}
\cas{3.}
\begin{calculs}
P(S\cap L) &= P(S)\times P_{S}(L) &&\justif{règle du chemin}\\
P(S\cap L) &= 0{,}7\times 0{,}2 = 0{,}14 &&\\
P(\overline{S}\cap L) &= P(L)-P(S\cap L) &&\justif{probabilités totales :
$P(L)=P(S\cap L)+P(\overline{S}\cap L)$}\\
P(\overline{S}\cap L) &= 0{,}23-0{,}14 = 0{,}09 &&
\end{calculs}
\cas{4.}
\begin{calculs}
P_{\overline{S}}(L) &= \dfrac{P(\overline{S}\cap L)}{P(\overline{S})} &&\justif{formule
des probabilités conditionnelles}\\
P_{\overline{S}}(L) &= \dfrac{0{,}09}{0{,}3} = \res{0{,}3} &&
\end{calculs}
\explique{$30\,\%$ des locations du week-end durent plus d'une heure.}{on interprète}
\cas{5.}
\begin{calculs}
P_{L}(\overline{S}) &= \dfrac{P(\overline{S}\cap L)}{P(L)} &&\justif{probabilité
inversée}\\
P_{L}(\overline{S}) &= \dfrac{0{,}09}{0{,}23} = \res{\dfrac{9}{23}} &&\justif{on
multiplie par $100$ ; $23$ est premier}
\end{calculs}
\cas{6.}
\explique{$P_{S}(L)=0{,}2$ et $P(L)=0{,}23$ : $P_{S}(L)\neq P(L)$, donc $S$ et $L$ ne sont
\res{pas indépendants}.}{autre façon : $P(S)\times P(L)=0{,}161\neq 0{,}14$}
\end{exo}

\begin{exo}[Un jeu à la fête foraine]
\cas{1.}\par\nobreak
\noindent\begin{minipage}[t]{0.4\linewidth}\vspace{0pt}\raggedright
Avec remise : à chaque tirage, $P(R)=\dfrac{4}{10}=0{,}4$ et $0{,}6$ pour une verte,
aux deux niveaux.
\end{minipage}\hfill
\begin{minipage}[t]{0.58\linewidth}\vspace{0pt}\centering
\begin{tikzpicture}
  \arbfeuilles{0{,}16}{0{,}24}{0{,}24}{0{,}36}
  \arbre{R_{1}}{R_{2}}{0{,}4}{0{,}6}{0{,}4}{0{,}6}{0{,}4}{0{,}6}
\end{tikzpicture}
\end{minipage}
\cas{2.}
\explique{$0$ ; $1$ ou $2$ boules rouges : on reçoit $0$ ; $3$ ou $6$\,€, moins les $2$\,€
payés. $X$ prend les valeurs $\res{-2}$ ; $\res{1}$ et $\res{4}$.}{gain algébrique}
\cas{3.}
\explique{$X=1$ : exactement une boule rouge, rouge puis verte ou verte puis
rouge.}{deux chemins}
\begin{calculs}
P(X=1) &= P(R_{1}\cap\overline{R_{2}})+P(\overline{R_{1}}\cap R_{2}) &&\justif{deux
chemins}\\
P(X=1) &= 0{,}4\times 0{,}6+0{,}6\times 0{,}4 &&\justif{règle du chemin}\\
P(X=1) &= 0{,}24+0{,}24 = 0{,}48 &&
\end{calculs}
\begin{center}
\renewcommand{\arraystretch}{1.4}
\begin{tabular}{|c|c|c|c|}
\hline
$x_{i}$ & $-2$ & $1$ & $4$\\\hline
$P(X=x_{i})$ & $\res{0{,}36}$ & $\res{0{,}48}$ & $\res{0{,}16}$\\\hline
\end{tabular}
\end{center}
\rem{$P(X=-2)=0{,}6\times 0{,}6$ et $P(X=4)=0{,}4\times 0{,}4$ ; la somme vaut $1$.}
\cas{4.}
\begin{calculs}
E(X) &= 0{,}36\times(-2)+0{,}48\times 1+0{,}16\times 4 &&\justif{somme des
$p_{i}\,x_{i}$}\\
E(X) &= -0{,}72+0{,}48+0{,}64 &&\\
E(X) &= \res{0{,}4} &&
\end{calculs}
\explique{Sur un grand nombre de parties, le joueur gagne en moyenne $0{,}40$\,€ par
partie : le jeu lui est favorable.}{$E(X)>0$}
\cas{5.}
\explique{Avec une mise de $m$ euros, les gains deviennent $-m$ ; $3-m$ ; $6-m$.}{mêmes
probabilités}
\begin{calculs}
E(X) &= 0{,}36\times(-m)+0{,}48\times(3-m)+0{,}16\times(6-m) &&\justif{même formule}\\
E(X) &= -0{,}36m+1{,}44-0{,}48m+0{,}96-0{,}16m &&\justif{on développe}\\
E(X) &= 2{,}4-m &&\justif{on réduit}
\end{calculs}
\explique{Le jeu est équitable pour une mise de \res{$2{,}40$\,€}.}{$2{,}4-m=0$}
\end{exo}

\end{document}
